\chapter
[\CO{[JUS]} Rechtliche Grundlagen]
{\CC{[JUS]}\\Rechtliche Grundlagen}
\label{chp:JUS}
% \AasRsN{RS~III}

\emph{Maintainer: Christof Koller}

\AuthNotes{
Changelog:
2009-02-20: GABS: Fertig für Kurs 2009/02

2009-02-25: EDIT Koller

2009-04-22: GABS: Abkürzungen "`Pat."', "`KA"', "`KH"'
bereinigt und gegen Langform ausgetauscht, wie besprochen.}

\setcounter{tocdepth}{10}
\mtcsetdepth{minitoc}{3}
\minitoc
\mtcsetdepth{minitoc}{\value{DefaultMinitocDepth}}


\clearpage
\section{Allgemeine Rechtsgrundlagen}
% TODO : Section: Layout-Anpassung
\TODO{GABS}{yyyy-mm-dd}{Subsubsection: Layout-Anpassung notwendig.}{Dieser Unterunterabschnitt entspricht nicht dem Standard-Layout!}

\subsection{Struktur der Österreichischen Rechtsordnung}
% TODO : Subsection: Layout-Anpassung
\TODO{GABS}{yyyy-mm-dd}{Subsubsection: Layout-Anpassung notwendig.}{Dieser Unterunterabschnitt entspricht nicht dem Standard-Layout!}

%\begin{tabular}{| p{2.2cm} | p{2.2cm} p{2.2cm} p{2.2cm} | p{2.2cm} p{2.2cm}
% |}\toprule

\Layout{0}{\TxBeschreibung}{
\begin{table}[H]
\Captionof{table}{Rechtsordnung}
\begin{tabularx}{\linewidth}[H]{X X|X|X||X|X||}

& \multicolumn{5}{c}{\EE{$\Swarrow$ EU-Recht $\Searrow$}}
\\\addlinespace

& \multicolumn{5}{c}{\EE{Struktur der Österreichischen Rechtsordnung}}
\\\addlinespace

{ }
& \multicolumn{3}{c}{\centering Öffentliches Recht}
& \multicolumn{2}{c}{\centering {Privatrecht} }
\\\hhline{~-----}

\multicolumn{1}{p{2cm}||}{Rechts\-gebiete }
& Strafrecht
& VwR
& VfR
& Zivilrecht
& Sonder\-privat\-rechte \\\hhline{~-----}

\multicolumn{1}{p{2cm}||}{Wichtigste Gesetzesbereiche} &
StGB, StPO &
Berufs\-ge\-setze,\newline KAKuG, UbG &
B-VG; StGG &
ABGB &
Arbeits\-rechts\-ge\-setze,
Unter\-nehmens\-ge\-setz\-buch\\\hhline{~-----}

\end{tabularx}
\end{table}
%\ncbarr[arrows=->]{EU}{OR}


Die österreichische Rechtsordnung ist ein zusammenhängendes Geflecht
bestehend aus zwei Hauptbereichen, dem öffentlichen Recht und dem
Privatrecht. Während das öffentliche Recht die Beziehung zwischen
Staat und Bürger regelt, bestimmt das Privatrecht die Verhältnisse
unter den Bürgern (= Privaten) selber:

\begin{description}
    \item[\T{Öffentliches Recht}:] Staat ${\leftrightarrow}$ Bürger
    \item[\T{Privatrecht}:] Bürger ${\leftrightarrow}$ Bürger
\end{description}

Durch den Beitritt zur \E{EU} wurde ein weiterer Gesetzgeber geschaffen,
der in beide Rechtsgebiete einwirken kann. Unabhängig davon lässt
sich jedes von beiden Rechts\-ge\-bieten in folgende
Bereiche unterteilen:
}{
\begin{itemize}
  \item Öffentliches Recht
  \begin{itemize}
    \item Strafrecht
    \item Verwaltungsrecht
    \item Verfassungsrecht
  \end{itemize}
  \item Privatrecht
  \begin{itemize}
    \item Zivilrecht
    \item Sonderprivatrechte
  \end{itemize}
\end{itemize}
}{}{}{}

\Layout{2}{Öffentliches Recht}
{
\begin{itemize}
    \item \T{Strafrecht}: Erfasst jene Delikte, die der
    Gesetzgeber als derart schwerwiegend erachtet hat, dass
    sie nicht nur (idR.) vom Staatsanwalt vor Gericht zu
    verfolgen sind,  sondern auch (in den meisten Fällen) mit
    Haftstrafen  sanktionierbar sind. Der Staatsanwalt
    vertritt hierbei  den Staat als Ankläger. Die eingehobene
    Geldstrafe  erhält der Staat, nicht der Verletzte.
    
    \item \T{Verwaltungsrecht}: Umfasst Regelungen die sich idR. nur an bestimmte Personenkreise richten und von Verwaltungsbehörden
    (nicht von Gerichten) bei Verstoß insofern sanktioniert werden, dass (vorrangig) Geldstrafen verhängt werden.
    
    \item \T{Verfassungsrecht}: Jener Teil der Rechtsordnung, der die grundlegenden Bestimmungen über Staatsgewalt, ihre Organe und
    ihre Verhältnisse zu den Staatsbürgern sowie
    Grundsätze der Rechtsordnung enthält. Er umfasst
    weiters Grundprinzipien und Grundrechte einer
    Rechtsordnung. 
    
\end{itemize}
}{}{}{}{}



\Layout{2}{Privatrecht}
{
\begin{itemize}
    \item \T{Zivilrecht}: Enthält jene allgemeinen Normen, die
    die Rechtsverhältnisse zwischen gleichgestellten Bürgern regeln.
    Bei Verstößen muss der Verletzte selber vor Gericht sein Recht
    einfordern (kein Staatsanwalt), wodurch er vom Schädiger
    Ersatzleistungen für den Schaden fordern kann (keine Strafe).
    
    \item \T{Sonderprivatrechte}: eigenständige Rechtsgebiete,
    die mittels selbstständigen Normen besondere Rechtsverhältnisse
    zwischen den Bürgern regeln. Zu den wichtigsten zählen das Arbeits-
    und das Unternehmensrecht.
    
\end{itemize}
}{}{}{}{}

\Layout{2}{Zugang zu den Gesetzestexten}
{
~
% POI 28

\noindent\includegraphics[width=\linewidth]{./images/JUS/RisHeader.jpg}

Die Gesetzestexte können im Internet vom Rechtsinformationssystem des Bundes
(RIS) unter der URL \url{http://www.ris.bka.gv.at/} abgerufen werden.

}{}{}{}{}

\subsection{Vertragsverhältnis im Rettungswesen}
\label{topic:vetrag-rettungswesen}
% TODO : Subsection: Layout-Anpassung
\TODO{GABS}{yyyy-mm-dd}{Subsubsection: Layout-Anpassung notwendig.}{Dieser Unterunterabschnitt entspricht nicht dem Standard-Layout!}

\Layout{0}{\TxBeschreibung}{
\begin{figure}[H]
%     \AuthNotes{KOLC: Graphik gehört rotiert.
%
%     GABS: ok, dauert aber noch ein bisserl. Wenn Inhalt
%     passt bitte etwas Geduld.}
\begin{center}
\includegraphics[width=8cm]{./images/recht_vertrag.png}
\end{center}
\Captionof{table}{Vertragsverhältnisse im Rettungsdienst}
\end{figure}
Das Rechtsverhältnis im Rettungswesen besteht grundsätzlich aus drei
Parteien: dem Patienten, der Rettungsorganisation
(Einrichtung) und dem (Not-)Arzt oder Sanitäter (Personal).
(In einer Krankenanstalt ist das Verhältnis ähnlich, mit dem Unterschied, dass an die Stelle der
Rettungsorganisation die Krankenanstalt und an jene des Sanitäters
der Krankenpfleger tritt.)

Als wichtigstes Verhältnis in dieser Dreiecksbeziehung ist das
\EE{Verhältnis zwischen Patient und
Rettungsorganisation} zu bezeichnen, da direkter
(Vertrags-)Partner zum Patient niemals der Sanitäter, sondern stets die von ihm vertretene
Organisation ist. Man nimmt nämlich an, dass ein Patient, wenn er eine
"`Rettung"' verständigt, nicht mit dem
einzelnen Einsatzpersonal in rechtlichen Kontakt treten möchte,
sondern mit der Gesamtheit einer Organisation. (Er bestellt ja nicht
dieses eine Fahrzeug mit jenem speziellen Personal.) Das
Rechtsverhältnis kann nun auf zwei verschiedene Arten zustande
kommen:

\begin{itemize}
  \item \T{Heilbehandlungsvertrag}: Der Patient selber oder sein
  zuständiger Vertreter	wünscht behandelt zu werden
     \item \EE{\T{Geschäftsführung ohne Auftrag} im Notfall}: der
     Patient befindet sich in einem derartigen Notfall, dass er sich nicht mehr
     äußern kann und dringend eine Behandlung benötigt (Bsp.:
     Bewusstlosigkeit)
\end{itemize}

Aus beiden Verhältnissen schuldet die Rettungsorganisation die
höchstmögliche Sorgfalt, aber niemals einen bestimmten
Behandlungserfolg. Weiters hat sie vorrangig nach dem
Willen des Patienten vorzugehen und wenn dieser fehlt, nach
seinem hypothetischen Willen.

Im \EE{Verhältnis zwischen Patient und
Personal} fungiert das Personal
bloß als Gehilfe zur Erfüllung der Pflichten seiner Organisation gegenüber
dem Patienten, wodurch man ihn auch als
\T{Erfüllungsgehilfen} bezeichnet. Er ist nicht
Vertragspartner des Patienten, muss aber mit der gleichen Sorgfalt, wie es
von der Organisation verlangt wird, arbeiten. Kommt es zu
Schädigungen des Patienten durch Handlungen des
Personals, kann der Geschädigte deliktisch gegen den
Erfüllungsgehilfen und vertraglich gegen die Organisation vorgehen. (IdR. wird der
Patient die Organisation verklagen, da ihm hierbei aufgrund
des Vertrages mehrere gesetzliche Erleichterungen zustehen.)

Das \EE{Verhältnis zwischen Rettungsorganisation
und Personal} lässt sich auf drei verschiedene Arten
begründen:

\begin{itemize}
    \item \E{Arbeitsvertrag}: Der (Not-)Arzt/Sanitäter ist bei
    der Organisation als Hauptamtlicher berufsmäßig beschäftigt.
    
    \item \E{Vereinsmitgliedschaft}: Der (Not-)Arzt/Sanitäter ist
    bei der Organisation als Vereinsmitglied und somit als
    Ehrenamtlicher/Freiwilliger tätig.
    
    \item \E{Zivildienst}: Der Sanitäter leistet seinen
    Zivildienst bei der Organisation.
    
\end{itemize}

Aus diesem Verhältnis heraus müssen manche der Pflichten, welche die
Organisation gegenüber dem Patienten treffen, auch
durch das Personal erfüllt werden. Gleichzeitig
treffen auch die Organisation und das Personal Rechte und
Pflicht im Umgang miteinander.
}{
\begin{itemize}
  \item Verhältnis zwischen Patient und Rettungsorgansiation
  \begin{itemize}
    \item Heilbehandlungsvertrag
    \item Geschäftsführung ohne Auftrag im Notfall
  \end{itemize}
  \item Verhältnis zwischen Patient und Personal
  \begin{itemize}
    \item Personal ist Erfüllungsgehilfe
  \end{itemize}
  \item Verhältnis zwischen Rettungsorganisation und
  Personal
  \begin{itemize}
    \item Arbeitsvertrag
    \item Vereinsmitgliedschaft
    \item Zivildienst
  \end{itemize}
\end{itemize}
}{}{}{}

\subsection{Strafbarkeitsprüfung}
\label{topic:strafbarkeitspruefung}
% TODO : Subsection: Layout-Anpassung
\TODO{GABS}{yyyy-mm-dd}{Subsubsection: Layout-Anpassung notwendig.}{Dieser Unterunterabschnitt entspricht nicht dem Standard-Layout!}

Überblick:
\begin{enumerate}
  \item Tatbestandserfüllung bzw. Schaden
  \item Kausalität
  \item Rechtswidrigkeit
  \begin{itemize}
    \item
    Rechtfertigungsgründe:
    \begin{itemize}
      \item Heilbehandlung
      \item Einwilligung
      \item Notwehr
      \item rechtfertigender Notstand
      \item Anhalterecht Privater
    \end{itemize}
  \end{itemize}
  \item Verschulden
  \begin{itemize}
    \item
    Entschuldigungsgründe:
    \begin{itemize}
      \item Notwehrüberschreitung im Affekt
      \item entschuldigender Notstand
      \item (Unzurechenbarkeit)
    \end{itemize}
  \end{itemize}
\end{enumerate}

Damit es zur Haftung für eine gesetzte Handlung kommt, müssen die
oben überblicksartig dargestellten und folgend näher beschriebenen
Voraussetzungen vollständig erfüllt sein. Mangelt es auch nur an
einer von ihnen, kommt es zu keiner Haftung. Am häufigsten kommen die
oben aufgezählten Rechtfertigungs- oder Entschuldigungsgründe ins
Spiel.

\Layout{2}{Tatbestandserfüllung}
{
Erfüllung der rechtlichen Beschreibung einer Lebenssituation auf die
eine Rechtsfolge/Strafe angeordnet wird.
}{}{}{}{}


\Layout{2}{Kausalität}
{

Die gesetzte Handlung war der Grund für den Eintritt der
Tatbestandserfüllung
}{}{}{}{}


\Layout{2}{Rechtswidrigkeit}
{
Die gesetzte Handlung war objektiv falsch ${\rightarrow}$ Urteil über
Tat ${\rightarrow}$ Verstoß gegen gesetzliche Verbote bzw. Gebote,
vertraglichen Bestimmungen oder die guten Sitten


\begin{description}
    \item[Rechtfertigungsgründe]  Der begangenen Tat kann kein Urteil gemacht werden, da sich ein
    maßgerechter Mensch in  derselben Situation gleich verhalten
    hätte
\end{description}
}{}{}{}{}


\Layout{2}{Verschulden}
{
Die gesetzte Handlung war subjektiv falsch ${\rightarrow}$ Urteil über
Täter

\begin{description}
    \item[Entschuldigungsgründe] Über den Täter kann kein Urteil gemacht werden, da er sich nach
    seinen persönlichen   Fähigkeiten und Kenntnissen in dieser
    Situation nicht anders verhalten hätte können
\end{description}



Die einzelnen Rechtfertigungs- und Entschuldigungsgründe werden in der
Ausarbeitung "`Ein langer Dienst"'
separat behandelt.
}{}{}{}{}





\clearpage

\section{Ein langer Dienst}
\label{topic:langer-dienst}
% TODO : Section: Layout-Anpassung
\TODO{GABS}{yyyy-mm-dd}{Subsubsection: Layout-Anpassung notwendig.}{Dieser Unterunterabschnitt entspricht nicht dem Standard-Layout!}

{\centering
\EE{{\sffamily\Large Eine Geschichte aus 17 Rechtsfällen}}
\par}

\Layout{2}{Zweck}
{
Behandlung rechtlicher Fragen anhand von Praxisbeispielen, die in einer
zusammenhängenden Geschichte erzählt werden. Dabei wird
insbesondere versucht auf die Besonderheiten des Rettungswesens Bezug
zu nehmen und die sich daraus ergebenden rechtlichen Probleme einfach
darzustellen. Die Darstellung beinhaltet aber nur einen Überblick
über die behandelten Rechtsgebiete sowie eine Konkretisierung auf die
am häufigsten gestellten Rechtsfragen.
}{}{}{}{}



\Layout{2}{Aufbau}
{
Jeder einzelne Rechtsfall wird in der Geschichte nach folgendem Aufbau
gestaltet:

\begin{enumerate}
    \item Sachverhalt: Schilderung der Geschichte
    \item Rechtsfrage: Erkennung des sich stellenden Problems
    \item Rechtslage: rechtliche Beurteilung des Problems
    \item Rechtsfolgen: drohende Konsequenzen und Ausweichmöglichkeiten
\end{enumerate}
}{}{}{}{}


\Layout{2}{Vorkommende Personen}
{
\begin{center}
\begin{tabular}{m{6.649cm}m{6.651cm}}
\textbf{\textcolor{red}{A}}lex, Fahrer

\textbf{\textcolor{red}{S}}epp, Transportführer

\textbf{\textcolor{red}{B}}ernd, Zivildiener

\textbf{\textcolor{red}{NA}}tascha, Notärztin

&
Hr. \E{Hauzu}, 1. Patient

Fr. \E{Uninformiert}, 2. Patientin

Fr. \E{Finito}, 3. Patientin

Hr. \E{Schmerz}, 4. Patient

Fr. \E{Aufgeregt}, Angehörige

Dr. \E{Wasbringst}, Aufnahmearzt

Fr. \E{Hilflos}, 5. Patientin

\\
\end{tabular}
\end{center}
}{}{}{}{}

\Layout{2}{Der Tag beginnt \ldots}
{%
Unsere Geschichte beginnt an einem warmen Frühjahrs\-morgen auf den die
drei Sanitäter Alex, Sepp und Bernd schon lange gehofft hatten, da
sie endlich nach dem langen Winter einen angenehm warmen Sonnentag
genießen wollen. Doch wie immer an einem solchen Freudentag, ruft
die Arbeit mit einem ganz normalen Sanitätsdienst und seinen
typischen Problemen.
}{}{}{}{}


\subsection{Gefahrenbekämpfung}

\subsubsection{Fall 1. Randalierender Patient}
\label{topic:unterbringung-jus}
% TODO : Subsubsection: Layout-Anpassung
\TODO{GABS}{yyyy-mm-dd}{Subsubsection: Layout-Anpassung notwendig.}{Dieser
Unterunterabschnitt entspricht nicht dem Standard-Layout!}


\Layout{2}{Sachverhalt}{
Die erste Ausfahrt beginnt natürlich direkt nach Dienstbeginn, ohne
dass wenigstens ein Kaffee getrunken werden hätte können. Der auf
unsere RTW-Mannschaft wartende Patient, Hr. Hauzu, hingegen hat die
ganze Nacht durchgemacht und ist schon weit\-läufig bekannt für seine
tobend aggressiven Wutausbrüche im Zuge seiner psychiatrischen
Erkrankung. Als Alex, Sepp und Bernd eintreffen droht er --
aus einiger Entfernung -- sich, oder einen anderen, zu
verletzen, falls man sich ihm auch nur nähern würde.
}{}{}{}{}


\subparagraph{Variante "`unmittelbare Bedrohung"':} Herr Hauzu
stürmt auf Alex zu, schmeißt ihn zu Boden und holt zum
Schlag aus.

\Layout{2}{Rechtsfrage}{Darf man Hr. Hauzu festhalten? Muss Hr. Hauzu eingewiesen werden?}{}{}{}{}


\Layout{2}{Rechtslage}{}{}{}{}{}

\EE{Eine Festhaltung \underbar{darf} erfolgen nach
folgenden Gesichtspunkten}:

\begin{itemize}
  \item Notwehr \Z{\S 3 Abs.~1 StGB}: Abwehr eines Angriffes
  auf ein eigenes oder fremdes Rechtsgut, wenn:
  \begin{itemize}
    \item gegenwärtig oder unmittelbar drohender Angriff,
    \item rechtswidrig,
    \item von Menschen gesetzt,
    \item auf ein notwehrfähiges Gut (= Leben, Gesundheit, Freiheit oder
    Vermögen) und
    \item das Gelindeste der zur verlässlichen Abwehr verfügbaren Mittel
    verwendet wird,
    \item welches nur gegen den
    Angreifer gerichtet werden darf.
  \end{itemize}
  \item \E{Notwehrüberschreitung im Affekt} \Z{{\S}~3
  Abs. 2 StGB}:
  \begin{itemize}
    \item Sollten Alex, Sepp und Bernd fälschlich einen
    Angriff angenommen haben, sind sie entschuldigt, wenn sie
    aus Furcht gehandelt haben und eine maßgerechte Person
    sich gleich verhalten hätte.
  \end{itemize}
  \item  \E{Anhalterecht Privater} \Z{\S 80 Abs. 2
  StPO}: Festhaltung bei Verdacht auf eine gerichtlich strafbare Handlung, welche im vorliegenden Fall
  die Gefahr einer Körperverletzung (KV) iSd. {\S}{\S}
  83\,ff StGB und die Sachbeschädigung iSd. {\S}\,125 StGB
  darstellt. Dabei darf nur auf angemessene Weise (Nötigung und leichte KV) angehalten werden, wenn
  unverzüglich für eine Meldung an die
  Polizei\MarginNo{\includegraphics[width=\linewidth]{./images/polizei_funkwagen.jpg}}
  gesorgt wird.
\end{itemize}

\EE{Eine Unterbringung \underbar{muss} erfolgen, nach
folgenden Gesichtspunkten}:

\begin{itemize}
  \item \Z{{\S}\,3} Unterbringungsgesetz (UbG):
  \begin{enumerate}
    \item psychiatrische Krankheit
    \item Selbst- und/oder Fremdgefährdung
    \item Verhältnismäßigkeitsgrundsatz
  \end{enumerate}
\end{itemize}
\begin{itemize}
  \item zu der Feststellung von a) und b) benötigt es das Gutachten eines
  \E{Amtsarztes}, welcher der Bezirksverwaltungsstelle (in Wien die MA\,15)
  unterstellt ist und über die Polizei zu verständigen ist. Der
  Amtsarzt darf die "`Zuweisung"'
  bis zur Krankenanstalt bestimmen, in der Krankenanstalt
  wird von 2 Fachärzten der Psychiatrie die eigentliche
  "`Einweisung"' verfügt. Während
  der Untersuchung und dem Transport darf Zwang angewendet werden,
  aber nur durch die gelindesten zur Erfolgsherbeiführung möglichen
  Mittel bei Beachtung der Mittel-Ziel-Relation (=\,Verhältnismäßigkeitsgrundsatz).
\end{itemize}

\Layout{2}{Rechtsfolge}{}{}{}{}{}
\Fazit{
Alex, Sepp und Bernd haben somit \E{grundsätzlich} das
Recht, eventuelle Gewaltausbrüche des Hr. Hauzu gegen sich
und andere abzuwehren. In der Situation, in der bloß aus einiger
Entfernung mit Gewalt gedroht wird, \E{ohne dass eine
unmittelbare Gefahr vorliegt}, ist
es \E{verhältnismäßig}, zurück zu treten, die Wohnung zu
verlassen und die Polizei zu verständigen. 

\subparagraph{Variante "`unmittelbare Bedrohung"':} Alex, Sepp
und Bernd können das verhältnismäßig gelindeste Zwangsmittel wählen um sich
gegen den unmittelbaren Angriff zu \E{schützen}. Sie sind
jedenfalls \E{verpflichtet}, unverzüglich die Polizei und
über diese den Amtsarzt zu verständigen. Trotzdem sollte Gewalt durch Sanitäter immer als letzte Option verwendet werden, auch wegen dem Risiko sich durch das Anhalten selbst zu verletzen. Daher der Grundsatz: Selbstschutz geht vor Fremdschutz (bzw. übertriebenen Heldenmut).
}
\subsubsection{Fall 2. Gefahr durch Hund}
% TODO : Subsubsection: Layout-Anpassung
\TODO{GABS}{yyyy-mm-dd}{Subsubsection: Layout-Anpassung notwendig.}{Dieser Unterunterabschnitt entspricht nicht dem Standard-Layout!}


\Layout{2}{Sachverhalt}{
Kaum hat Hr. Hauzu gehört, dass die Polizei und der Amtsarzt
verständigt wurden, stürmt er zu einem kleinen Nebenzimmer in
seiner Gemeindebauwohnung, öffnet die Tür und lässt einen
"`kleinen"' (65\,kg schweren)
Rottweiler-Kampfhund heraus. Mit fletschenden Zähnen nähert
er sich den Sanitätern und droht jeden Moment anzugreifen.
}
{}{}{}{}



\Layout{2}{Rechtsfrage}{Was darf gegen den Hund getan werden?}{}{}{}{}


\Layout{2}{Rechtslage}{}{}{}{}{} ~


\begin{table}[H]\begin{WidePar}
\Captionof{table}{Rechtslage: Gegenüberstellung rechtfertigender vs.
entschuldigender Notstand.}

\begin{tabularx}{\linewidth}{XX}
\EE{Rechtfertigender Notstand}
&
\EE{Entschuldigender Notstand}
\\\midrule
\begin{itemize}
    \item in keinem Gesetz aufgezeichnet; wird aber einheitlich
    angenommen
    \item anzuwenden wenn die Gefahr auch tatsächlich
    (objektiv) vorliegt
\end{itemize}
&
\begin{itemize}
    \item {\S}\,10 StGB; {\S}\,1306a ABGB
    \item anzuwenden wenn man die Gefahr bloß subjektiv
    annimmt
\end{itemize}
\\\midrule

\EE{Notstandssituation}: Gefahr ist:

\begin{itemize}
    \item  kein menschlicher Angriff
    \item  gegenwärtig oder unmittelbar drohend
    \item rechtswidrig
    \item  auf beliebiges Rechtsgut gerichtet
    \item  droht bedeutender Nachteil
\end{itemize}

\EE{Notstandshandlung}:

\begin{itemize}
    \item  Abwehr als schonendstes Mittel
    \item  Zurechnungsprinzip (=~Abwehr geht in Sphäre des
    Gefahrenverursachers)
    \item  geringes Risiko bei hoher Rettungschance 
\end{itemize}
&
\EE{Notstandssituation}: Gefahr ist:

\begin{itemize}
  \item  kein menschlicher Angriff
  \item  gegenwärtig oder unmittelbar drohend
  \item rechtswidrig
  \item  auf beliebiges Rechtsgut gerichtet
  \item  droht bedeutender Nachteil
\end{itemize}

\EE{Notstandshandlung}:

\begin{itemize}
  \item  Abwehr als schonendstes Mittel
  \item  Zurechnungsprinzip (= Abwehr geht in Sphäre des
  Gefahrenverursachers)
  \item  Drohender Nachteil darf geringfügig kleiner als das
  Ergebnis der Abwehrhandlung sein.
\end{itemize}
\\\bottomrule
\end{tabularx}\end{WidePar}
\end{table}


\Layout{2}{Rechtsfolge}{}{}{}{}{}

\Fazit{
Alex, Sepp und Bernd dürfen sich mit den schonendsten Mitteln
verteidigen, um ihr eigenes Leben oder ihre Gesundheit zu schützen.}

\begin{quotation}
Praxistipp: Schon vor Betreten der Wohnung sollte ein Hund von seinem
Besitzer in einen  separaten Raum weggesperrt werden.
\end{quotation}

\subsection{Heilbehandlung }

\subsubsection{Fall 3. Blutzuckermessung =  Heilbehandlung }
% TODO : Subsubsection: Layout-Anpassung
\TODO{GABS}{yyyy-mm-dd}{Subsubsection: Layout-Anpassung notwendig.}{Dieser Unterunterabschnitt entspricht nicht dem Standard-Layout!}



\Layout{2}{Sachverhalt}{
Durch die rechtzeitig eintreffende Polizei, kann die Gefahr durch den
Hund beseitigt und Hr. Hauzu festgehaltenen werden. Trotz der Zuweisung
durch den etwas später eingetroffenen Amtsarzt, möchte Bernd,
gewissenhaft wie er ist, einen vollen Neuro-Check mit einem
Blutzuckertest durchführen.}{}{}{}{}



\Layout{2}{Rechtsfrage}{
Wann liegt eine  Heilbehandlung  (= HB) vor? Darf der Blutzuckertest
durchgeführt werden?
}{}{}{}{}


\Layout{2}{Rechtslage}{

\T{ Heilbehandlung }:
Wird gewonnen aus allgemeinen Rechtsgrundsätzen und der systematischen
Betrachtung der gesamten Rechtsordnung.
}{}{}{}{}

\subparagraph{Voraussetzung:}

\begin{itemize}
    \item medizinisch indiziert (= Heilbehandlung ist zur Diagnose oder Linderung von Erkrankungen und/oder Verletzungen notwendig)
    \item lege artis durchgeführt (= Heilbehandlung entspricht dem geltenden Stand der Wissenschaft)
    \item durchzuführen nach Aufklärung des Patienten
    (siehe \prettyref{topic:jus-fall-aufklaerung}).
\end{itemize}

\Fazit{
\noindent{}${\rightarrow}$ das Vorliegen einer Heilbehandlung wird als
Rechtfertigung für alle Einwirkungen am Patienten gewertet

\noindent{}${\rightarrow}$ wer die Voraussetzungen der Heilbehandlung nicht
einhält macht sich strafbar nach:
}

\begin{itemize}
    \item Strafrecht:
    
    \begin{itemize}
        \item a) \EE{Körperverletzung} - {\S}{\S}\,83 ff StGB: Sonderfall der HB-Voraussetzungen: es benötigt nur med. Indikation und lege artis Durchführung; die Aufklärung ist hierfür irrelevant
        
        \item b) \EE{Eigenmächtige  Heilbehandlung } - {\S}\,110 StGB: siehe ad 4.
        
    \end{itemize}
    \item Zivilrecht:
    \begin{itemize}
        \item a) \EE{Schadenersatzansprüche} - {\S}{\S}
        1293 ff ABGB: die Heilbehandlung muss alle drei
        Voraussetzungen erfüllen (auch Aufklärung);
        Umfang des Schadenersatzes siehe \prettyref{topic:jus-fall-fortbildung}.
    \end{itemize}
    
\end{itemize}

\subparagraph{Blutzuckermessung}

Die Messung des Blutzucker mit einer Lanzette war früher umstritten,
da manche Rettungsorganisationen davon ausgingen, dass es sich um einen
nicht dem Sanitäter erlaubten Eingriff am menschlichen Körper
handelt. Dieser Streit wurde durch eine Novellierung des SanG gelöst,
indem man die Blutzuckermessung für den Sanitäter ausdrücklich
zuließ.

\Layout{2}{Rechtsfolge}{}{}{}{}{}
\Fazit{
Da ein kompletter Neuro-Check bei einem neurologisch
auffälligen Patienten medizinisch indiziert und
gesetzlich erlaubt ist, darf jeder der drei Sanitäter eine solche Maßnahme lege artis durchführen. Eine
Aufklärung ist aufgrund der mangelnden Einsichts- und
Urteilsfähigkeit nicht möglich.}

\subsection{Selbstbestimmungsrecht}

\subsubsection{Fall 4. Aufklärung und Einwilligung}
\label{topic:jus-fall-aufklaerung}
% TODO : Subsubsection: Layout-Anpassung
\TODO{GABS}{yyyy-mm-dd}{Subsubsection: Layout-Anpassung notwendig.}{Dieser Unterunterabschnitt entspricht nicht dem Standard-Layout!}

\Layout{2}{Sachverhalt}{
Nachdem nun Hr. Hauzu gebändigt, zur Gänze untersucht und
eingewiesen wurde, freut sich die RTW-Mannschaft schon auf den
wohlverdienten Kaffee am Stützpunkt. Doch wie es kommen musste,
können Alex, Sepp und Bernd nicht einrücken, da sie einen
Folgetransport zugewiesen bekommen. Die vorgefundene
Patientin Fr. Uninformiert ist eine in der 30. Wo.
schwangere Mutter, die erst letzte Woche ihren 18. Geburtstag feierte. Grund für ihren Anruf sind starke
Kopfschmerzen und leichte Seh- und Höreinschränkungen. Bei der
weiteren Untersuchung durch Sepp werden geschwollene Knöchel, ein
Blutdruck von 140\,/\,90\,mm\,Hg und nach beiläufiger Aussage der
Patientin starker Harndrang mit süßlichem Geruch
diagnostiziert.

Schnell ist Sepp klar, dass es sich um eine \index{EPH-Gestose} \E{EPH-Gestose}
(\ref{topic:eph-gestose}) handelt, die zur Eklampsie führen kann und einen raschen, aber
möglichst reizarmen Transport auf die Gynäkologie benötigt. Da er aber wenig von jungen
Müttern hält, erachtet er es auch nicht für notwendig Fr.
Uninformiert über die nächsten Geschehnisse zu informieren, sondern
ordnet nur harsch Befehle an, die von Fr. Uninformiert
"`gefälligst"' durchzuführen sind. 
}{}{}{}{}

\Layout{2}{Rechtsfrage}{
Besteht eine Aufklärungspflicht und was umfasst sie? Muss eine
Einwilligung erfolgen und wenn ja, in welcher Form?
}{}{}{}{}


\Layout{2}{Rechtslage}{%
%\begin{multicols}{2}
\subparagraph{\EE{Aufklärungspflicht}:}

\begin{itemize}
    \item Jeder Patient ist aufzuklären,
    \item vor Setzung der Maßnahme,
    \item wenn er einsichts- und urteilsfähig ist.
\end{itemize}

\Fazit{Seit dem 1. Dezember 2009 wurde das \EE{Recht auf
Aufklärung} durch das Inkrafttreten des Artikel~3 Abs.\,2
lit.\,a der Grundrechtecharta der \Abk{EU} zu einem
eigenen \EE{Grundrecht}.} 

\subparagraph{\EE{Einsichts- und Urteilsfähigkeit}:}

\begin{itemize}
  \item grundsätzlich: bei jedem volljährigen Patienten gegeben
  \item \E{fehlend bei}:
  \begin{enumerate}
    \item grundsätzlich unmündigen Minderjährigen (\VonBis{0}{14} LJ.), wobei es
    vorrangig auf die Einsichts- und Urteilsfähigkeit ankommt
    \item bei mündig Minderjährigen (\VonBis{14.}{18.} LJ.) vorhanden, solange
    entsprechende Einsichts- und Urteilsfähigkeit nicht fehlt
    \item psychisch erkrankter Patient mit fehlender Einsichts- und Urteilsfähigkeit
    \item Patient im Schockzustand \AuthNotes{GABS: CAVE
    Verwechslung mit "`echtem"' Schock -- anders Formulieren?}
    \item stark berauschter Patient (ab ca. 2,5{\textperthousand})
    \item Bewusstlosen
  \end{enumerate}
  \item[] ad 1) \& 2): {\S}\,146c ABGB:
  \begin{itemize}
    \item grundsätzlich entscheidend ist die vom Einsatzpersonal getroffene
    Beurteilung der Einsichts- und Urteilsfähigkeit des Minderjährigen
    \item nur im Zweifel ist erst der mündige Minderjährige selbstständig
    entscheidungsbefugt (es benötigt keine Zustimmung der Eltern)
    \item nur bei schwerwiegenden Angelegenheiten muss ein Elternteil zustimmen
    ${\rightarrow}$ stimmt ein Elternteil zu, der andere verweigert aber,
    so ist  entweder jener Teil entscheidungsbefugt der die größere
    Rolle in der Erziehung des Patienten spielt (im Zweifel
    die Mutter) oder das \index{Pflegschaftsgericht}Pflegschaftsgericht ist anzurufen
    \item lehnen beide ab, obwohl es sich um eine notwendige Heilbehandlung handelt, muss das
    Pflegschaftsgericht angerufen werden
  \end{itemize}
\end{itemize}

\subparagraph{\EE{Umfang der Aufklärung}:}

\begin{itemize}
  \item Grundsatz: Die Aufklärung hat in jenem Rahmen zu erfolgen, der einer
  maßgerechten Figur unter Berücksichtigung der individuellen
  Besonderheiten des Patienten und der Situation gerecht
  wird.
  \item Detail:
  \begin{itemize}
    \item Umfang und Tragweite der  Heilbehandlung
    \item Folgen, Gefahren und Nebenwirkungen der
    Heilbehandlung
    \item Umstände die für verständigen Patienten ins
    Gewicht fallen würden
    \item umgekehrte Proportion zur Dringlichkeit der
    Heilbehandlung
  \end{itemize}
\end{itemize}

\subparagraph{\EE{Unterlassen der Aufklärung}:}
Unterbleibt die Aufklärung, wird angenommen, dass dem
Patient die Tragweite der Maßnahmen nicht bewusst ist und er daher nicht in die
Heilbehandlung einwilligt.

\subparagraph{\EE{Besonderheit \T{Revers}}:}\index{Revers}
\label{topic:revers}

\begin{itemize}
  \item Was ist der \T{Revers}?
  \begin{itemize}
    \item
    Der Revers ist eine schriftliche Dokumentierung über die Aussage des
    Patient, eine vom Einsatz\-personal empfohlenen Maßnahme
    abzulehnen. Die Verweigerung kann auch vom zuständigen Vertreter des
    Patienten abgegeben werden.
  \end{itemize}
  \item Welche Voraussetzungen hat der Revers?
  \begin{itemize}
    \item auszufüllen nach umfassender Aufklärung des Patienten
    \item Schriftlichkeit (Transportschein, Reversschein oder beigelegtes
    Dokument)
    \item Unterschrift des Patienten (oder seines zuständigen Vertreters)
  \end{itemize}
  \item Wann und warum ist der Revers notwendig?
  \begin{itemize}
    \item wenn das Einsatzpersonal eine Maßnahme für notwendig erachtet,
    der Patienten  (oder sein zuständiger Vertreter) diese Maßnahme
    aber verweigert
    \item zur Einhaltung der gesetzlichen Dokumentationspflicht und zur Abwendung
    von  klagsweise geltendmachbaren Ansprüchen des Patienten
  \end{itemize}
\end{itemize}
%\end{multicols}

\subparagraph{\EE{Besonderheit
\T{Gleichheitsgrundsatz}}:}\index{Gleichheitsgrundsatz}
\Z{(Art.~7 B-VG, Art.~2 StGG, Art.~20 Z.~1 EGC)}

\begin{quote}
\emph{Alle Patienten sind gleich zu behandeln, unabhängig ihrer
Religion, Hautfarbe, ethnische Zugehörigkeit,
geschlechtliche Orientierung oder ihres Geschlechts oder
sonstiger unsachlicher Unterscheidung.} 
\end{quote}
}{}{}{}{}


\paragraph{Rechtsfolgen}

\Fazit{
Aufgrund der unterlassenen Aufklärung machte sich Sepp strafbar nach
folgenden Bestimmungen:

\begin{itemize}
  \item \EE{Strafrecht}:
  \begin{itemize}
    \item Eigenmächtige  Heilbehandlung  ({\S}\,110 StGB):
    \begin{itemize}
      \item nur bei "`Gefahr in Verzug"'\marginlabel{Gefahr in Verzug: \E{unmittelbar} drohende oder
      gegenwärtige Gefahr, die ein Abwarten nicht zulässt bzw.
      unmittelbares Handeln verlangt} benötigt es
      keine Aufklärung bzw. Einwilligung und es dürfen alle nach Ansicht
      der Einsatzkräfte direkt notwendigen Maßnahmen gesetzt werden
    \end{itemize}
    \item Körperverletzung \Z{({\S}{\S}\,83 ff StGB)}:
    \begin{itemize}
      \item keine Strafbarkeit, solange die Heilbehandlung medizinische indiziert und lege artis
      durchgeführt wurde
    \end{itemize}
  \end{itemize}
  \item \EE{Zivilrecht}:
  \begin{itemize}
    \item
    Schadenersatz wegen Körperverletzung \Z{({\S}{\S}\,1293 ff ABGB)}:
    \begin{itemize}
      \item {für das Zivilrecht ist die Einwilligung aufgrund einer Aufklärung
      notwendige Voraussetzung für die rechtfertigende Heilbehandlung ${\rightarrow}$
      mangelt es an ihr, kann Schadenersatz verlangt werden}
    \end{itemize}
  \end{itemize}
\end{itemize}
}

\subsubsection{Fall 5. Patientenverfügung}
\label{topic:patientenverfuegung}
% TODO : Subsubsection: Layout-Anpassung
\TODO{GABS}{yyyy-mm-dd}{Subsubsection: Layout-Anpassung notwendig.}{Dieser Unterunterabschnitt entspricht nicht dem Standard-Layout!}
\Layout{2}{Sachverhalt}{
Nach Ablieferung der Fr. Uninformiert drängt Alex seine Kollegen zu
einer schnellen Abfahrt, um endlich seinen Kaffee zu bekommen. Doch da
macht ihm die Leitstelle schon wieder einen Strich durch die Rechnung,
als sie das Team zu einem weiteren BO beruft.

Dieser liegt in einem kleinen Pflegeheim, indem eine aufgeregte
Heimpflegerin auf Alex, Sepp und Bernd zustürmt und sie zu der 70
jährigen, regungslosen Patientin, Fr. Finito, bringt.
Gleichzeitig hält die Heimpflegerin den Sanitätern eine Patientenverfügung vor, die
Fr. Finito vor über 5 Jahren auf einer Serviette unterschrieben hat
und nach der "`bei ausbleibenden Kreislauf keine
Wiederbelebungsmaßnahmen gesetzt werden"'
dürfen. Deshalb hatte die Pflegerin es bis jetzt auch unterlassen mit
einer Reanimation zu beginnen. Da sie aber ebenso nicht weiß, ob
die Verfügung bindend ist, rief sie die Rettung und setzte sich
wartend neben die Patientin.
}{}{}{}{}

\Layout{2}{Rechtsfrage}{}{}{}{}{}

Was ist eine Patientenverfügung? Ab wann und wie weit ist sie bindend?

\Layout{2}{Rechtslage}{}{}{}{}{}

%\begin{multicols}{2}
\subparagraph{\EE{Allgemeines zur Patientenverfügung}:}

\begin{itemize}
  \item Willenserklärung einer einsicht-, urteils- und
  äußerungsfähigen Person
  \item erst von Bedeutung, wenn der Patient seinen Willen nicht mehr ausdrücken
  kann
  \item nur Ablehnung einer bestimmten Heilbehandlung kann Inhalt der Verfügung sein
  \item unbeachtlicher Inhalt:
  \begin{itemize}
    \item Ausschluss einer Grundversorgung
    \item med. nicht indizierte Behandlung
  \end{itemize}
\end{itemize}

\subparagraph{zwei Arten der Patientenverfügung:}

\begin{enumerate}
  \item \T{Verbindliche Patientenverfügung}:
  \begin{itemize}
    \item schriftlich vor Notar, Rechtsanwalt oder rechtskundigen
    Patientenvertreter
    \item eigenhändig unterschrieben und datiert
    \item aufgrund einer umfassenden Aufklärungen durch einen Arzt
    \item Inhalt muss eine konkrete Heilbehandlung ablehnen ${\rightarrow}$ allgemeine oder
    widersprüchliche Formulierungen führen zur Unwirksamkeit
    \item Wirkung nur für 5 Jahre
    \item Eine Verfügung die alle Elemente enthält ist absolut bindend und
    darf nicht umgangen werden
  \end{itemize}
  \item \T{Beachtliche Patientenverfügung}:
  \begin{itemize}
    \item
    \EE{grundsätzlich}: alle Verfügungen, die nicht
    die obigen Voraussetzungen erfüllen
    \item
    \EE{Umfang der Beachtlichkeit}: bloß relative
    Beachtlichkeit:
    \begin{itemize}
      \item
      Verfügung soll in Behandlungsentscheidungen berücksichtigt werden,
      aber
      \item
      das Wohl des Patienten bleibt oberstes Gebot
      \item
      daher kann man sich bei besonderer Begründung über die Verfügung
      hinwegsetzen
    \end{itemize}
  \end{itemize}
\end{enumerate}

\subparagraph{Absolute Unwirksamkeit/Unbeachtlichkeit bei:}

\begin{itemize}
  \item zu allgemeine oder widersprüchliche Formulierungen
  \item rechtswidriger Inhalt
  \item erkennbare Willensmängel des Patienten bei Abgabe der
  Verfügung (nicht ernst gemeint, Irrtum, List, Drohung oder Täuschung)
\end{itemize}


\subparagraph{Suche nach Patientenverfügung:}
Bei Notfallversorgungen muss nicht nach einer Verfügung gesucht
werden, wenn die Suche das Leben oder die Gesundheit des Patienten ernstlich
gefährden würde. Daher wird für das Rettungswesen idR. keine
Verpflichtung zur Suche, sondern zur Behandlung bestehen.\ZFootnote{{\S}\,12 PatVG}

%\end{multicols}

\Layout{2}{Rechtsfolge}{}{}{}{}{}

\Fazit{
Die Patientenverfügung gilt mit über 5 Jahren, der fehlenden
ärztlichen Aufklärung und dem mangelnden Beiziehen eines
Rechtsverständigen nicht mehr als verbindlich und wäre somit als
bloß beachtlich einzustufen.


Doch aufgrund des zu allgemeinen Inhalts, der Beiläufigkeit der
Verfügung (Serviette als Unterlage) und wegen des hohen Alters der
Verfügung kann kein konkreter und aktueller Wille der
Patientin erkannt werden, was dazu führt, dass das Wohl
der Patientin das oberste Gebot bleibt.

Mit einer Reanimation ist zu beginnen.

}

\subsection{Rechte und Pflichten des Sanitäters}

\subsubsection{Fall 6. Unterlassung der Notarztnachforderung}
% TODO : Subsubsection: Layout-Anpassung
\TODO{GABS}{yyyy-mm-dd}{Subsubsection: Layout-Anpassung notwendig.}{Dieser Unterunterabschnitt entspricht nicht dem Standard-Layout!}

\Layout{2}{Sachverhalt}{Endlich am Stützpunkt angekommen, ist Alex auch schon der Erste der
sich zum Kaffeeautomaten begibt. Entkräftet durch die einsetzende
"`Hypokoffeinomie"' kramt er nach
Kleingeld und kaum als er es einwurfbereit in Händen hält, tönt
die nächste Ausfahrt durch die Lautsprecher, mit dem Hinweis einer
dringenden Berufung. Alex geht aber in seiner typisch ruhigen Manie,
glücklich durch die enteral\footnote{Enteral: den Darm betreffend.} erlangte
Sedierung, zum Fahrzeug, mit dem Spruch auf den Lippen: "`immerhin konn des
Wagl{\textquoteright} eh net ohne mie foahrn!"'}{}{}{}{}



\begin{wrapfigure}{r}{5cm}
\includegraphics[width=\linewidth]{./images/teaser/nef-passat-kontur.jpg}
\end{wrapfigure}Nach nur kurzer Anfahrtszeit begegnet Alex, Sepp und Bernd im ersten
Stock eines kleinen Reihenhaus der 55 jährige, adipöse Hr. Schmerz
mit Schmerzen hinter der Brust, die in den linken Arm ausstrahlen und
ihm selbst in Ruhelage mit Atemnot und kalten Schweißausbrüchen
plagen. Für Sepp ist sofort klar, dass es sich um einen
MCI\footnote{\E{MCI}: Myokardinfarkt: Herzinfarkt.} handelt.
Er möchte schleunigst einen NEF nachfordern, doch Alex
meint, dass er keine Lust hätte schon wieder auf
einen "`Boder\footnote{"Bader": früher
Kundiger der Heilbäder, heute umgangssprachlich für
Arzt}"' zu warten und der Patient auch
selbst runter zum Fahrzeug gehen sollte, da er es im Kreuz
hätte und sowieso noch eine ca. 20 Minuten lange
Anfahrtszeit zur Krankenanstalt durch den schon einsetzende Morgenstau
bewältigen müsste.

\Layout{2}{Rechtsfrage}{%
Besteht die Pflicht für Sanitäter zur Nachforderung eines Notarztes?}
{}{}{}{}

\Layout{2}{Rechtslage}{

\EE{Unterlassene Hilfeleistung} - {\S}\,95 StGB

\begin{itemize}
    \item 
    jeder Durchschnittsmensch ist zur Durchführung von
    \E{ihm zumutbaren Hilfeleistungen} verpflichtet. Die
    Zumutbarkeit bestimmt sich dabei nach dem individuellen Stand an
    Ausbildung, Fähigkeiten und Kenntnissen. Aber schon zu den Grundlagen
    einer Ersten Hilfe, die von jedem Menschen mit Besuch eines
    Führerscheinkurses erwartet werden kann, gehört das Absetzen eines
    Notrufes zur  Anforderung von geeigneten Rettungskräften.
\end{itemize}

\EE{{\S}\,4 SanG}

\begin{itemize}
    
    \item 
    ein jeder Sanitäter, unabhängig von seinem Ausbildungsstand, ist zur
    Durchführung von lebensrettenden Sofortmaßnahmen verpflichtet,
    was gleichzeitig beinhaltet, dass er bei Notwendigkeit eine
    unverzügliche Anforderung des Notarztes oder sonstigen Arztes zu
    veranlassen hat.
    
\end{itemize}

\TODO{GABS}{2010-05-28}{Unterlassene Hilfeleistung: Abgrenzung von
Unterlassung von schädlichen Handlungen.}{
In der Diskussion wird oft jede Unterlassung mit unterlassener Hilfeleistung
gleichgesetzt. Es sollte ein gegenteiliger Hinweis erfolgen, zB.

"`Anmerkung: Die Unterlassung von schädlichen Handlungen, d.\,h. Handlungen die
\E{nicht} dem Patientenwohl dienen, stellt \E{keine} unterlassene Hilfeleistung
dar. Schädliche Dinge zu unterlassen ist genauso wichtig wie nützliche Dinge zu
tun."' 

Die Beurteilung, welche Maßnahme schädlich ist, ist oft stark vom Einzelfall
abhängig. }


}{}{}{}{}

\Layout{2}{Rechtsfolge}{}{}{}{}{}

\Fazit{
Würde Alex die zweckmäßige Nachforderung eines Notarztes unterlassen,
macht er sich gleichzeitig nach strafrechtlichen ({\S}\,95 StGB) und
verwaltungsrechtlichen Bestimmungen ({\S}\,9 in
Verbindung mit {\S}\,53 SanG) strafbar. }

% \clearpage
\subsubsection{Fall 7. Sanitäterkompetenzen}
\label{topic:kompetenzen-sanitaeter}
% TODO : Subsubsection: Layout-Anpassung
\TODO{GABS}{yyyy-mm-dd}{Subsubsection: Layout-Anpassung notwendig.}{Dieser Unterunterabschnitt entspricht nicht dem Standard-Layout!}

\Layout{60}{Sachverhalt}{Nach langem hin und her, können Sepp und Bernd
gemeinsam den Alex überreden doch einen Notarzt nachzufordern und Sepp, der gerade frisch
gebackener NFS geworden ist, sieht endlich seine Möglichkeit
umfangreich am Hr. Schmerz zu werken. Dabei möchte er neben der
Verabreichung von Nitrolingual-Spray auch gleich einen venösen Zugang
legen, da er das sowieso als ehemaliger Fleischhauer besser als jeder
Arzt kann.}
{}
{./images/venflon_in_situ.jpg}
{Ein blutiger Venflon vom Fleischhauer?}
{Gabriel.}



\Layout{2}{Rechtsfrage}{Was sind die \T{Kompetenzen} der einzelnen Sanitäterstufen
(RS, NFS, NKA, NKV, NKI)? Darf man die Kompetenzgrenzen (bei \E{Gefahr in
Verzug}) überschreiten?}{}{}{}{}



\Layout{2}{Rechtslage}{%
Die \EE{Kompetenzstufen} sind im SanG\ZFootnote{{\S}{\S}~\VonBis{9}{12} SanG}
geregelt. \prettyref{tab:sang-kompetenzstufen} gibt eine Übersicht:


% Freigabe Tabelle KOLC 2010-08-25
\Layout{57}{}{}{}{
\begin{tabulary}{\linewidth}{LL}
\TabHeading{Kompetenzstufe}
&
\TabHeading{Beschreibung}
\\\addlinespace\toprule\addlinespace
\TabSub{RS}

\E{Rettungssanitäter}
&
\begin{minipage}[t]{\linewidth}
\begin{itemize}
  \item Selbstständige und eigenverantwortliche Versorgung und Betreuung
    des Patienten
    \item Übernahme und Übergabe des Patienten im Rahmen
    eines Transportes
    \item Hilfestellung bei Akutsituationen; Verabreichung von \ce{O2}
    \item Durchführung von lebensrettenden Sofortmaßnahmen
    \item Durchführung von Sondertransporten
    \item Blutzuckermessung
\end{itemize}
\end{minipage}


\\\addlinespace\toprule\addlinespace
\TabSub{NFS}

\E{Notfallsanitäter}
&
\begin{minipage}[t]{\linewidth}
\begin{itemize}
  \item Tätigkeiten des RS
  \item Unterstützung des Arztes; einschließlich der
  Betreuung und des sanitätsdienstlichen Transports von \E{Notfallpatienten}
  \item Verabreichung von Medikamenten der der \EE{Arzneimittelliste 1}:
    Wird vom leitenden Arzt der Rettungsorganisation
    festgelegt. Die Verabreichung kann
    durchgeführt werden, ohne dass ein Notarzt informiert werden muss, oder es
    ein solcher anordnet.
    \item Eigenverantwortliche Betreuung von berufsspezifischen
    Medizinprodukten und Arzneimitteln
    \item Mitarbeit in Forschung
\end{itemize}
\end{minipage}

\\\addlinespace\midrule\addlinespace
\TabSub{NKA}

Notfallsanitäter mit \E{allgemeiner Notkompetenz Arzneimittellehre}
&
\begin{minipage}[t]{\linewidth}
\begin{itemize}
  \item Zusätzlich zu den oben genannten Kompetenzen darf der NKA jene
Medikamente der Arzneimittelliste 2 verabreichen, die ohne venösen
Zugang verabreicht werden können.
\item Die \EE{Arzneimittelliste 2} wird vom leitenden Arzt der Rettungsorganisation
festgelegt. Die Medikamente dürfen auf Anweisung eines anwesenden Arztes (jeder mit
jus practicandi, nicht unbedingt Notarzt) \E{oder} nach Verständigung eines
Notarztes verabreicht werden.
\end{itemize}
\end{minipage}
 \\\addlinespace\midrule\addlinespace
\TabSub{NKV}

Notfallsanitäter mit \E{allgemeiner Notkompetenz Venenzugang und Infusion}
& 
\begin{minipage}[t]{\linewidth}
\begin{itemize}
  \item Zusätzlich zu den oben genannten
  Kompetenzen darf der NKV jene Medikamente der Arzneimittelliste 2
  verabreichen, die über venösen Zugang verabreicht werden können.
\end{itemize}
\end{minipage}

\\\addlinespace\toprule\addlinespace
\TabSub{NKI}

Notfallsanitäter mit \emph{besonderer Notkompetenz Beatmung und Intubation} 
&
\begin{minipage}[t]{\linewidth}
\begin{itemize}
  \item Durchführung einer endotrachealen Intubation samt Beatmung, aber ohne
     Prämedikation. 
     \item Die NKI muss alle 2 Jahre rezertifiziert werden.
     \item Durchführung der Intubation auf Anweisung eines anwesenden Arztes (jeder mit jus
practicandi, nicht unbedingt Notarzt) \E{oder} nach Verständigung eines Notarztes.
\end{itemize}
\end{minipage}




\\\addlinespace\bottomrule
\end{tabulary}
}
{Kompetenzstufen von Sanitätern nach dem SanG ({\S}{\S}~\VonBis{9}{12})\label{tab:sang-kompetenzstufen}}
{} 


% In Tabelle übernommen und Freigabe von KOLC, zur Löschung vorgemerkt -- 2010-08-25
% 
% \begin{labeling}[~:]{XXXX}
%   \item[\T{RS}] \marginlabel{{\S}\,9 SanG}
%   \begin{itemize}
%     \item Z 1: selbstständige und eigenverantwortliche Versorgung und Betreuung
%     des Patienten
%     \item Z 2: Übernahme und Übergabe des Patienten im Rahmen
%     eines Transportes
%     \item Z 3: Hilfestellung bei Akutsituationen; Verabreichung von \ce{O2}
%     \item Z 4: Durchführung von lebensrettenden Sofortmaßnahmen
%     \item Z 5: Durchführung von Sondertransporten
%     \item Z 6: Blutzuckermessung
%   \end{itemize}
%   \item[\T{NFS}] \marginlabel{{\S}\,10 SanG}
%   \begin{itemize}
%     \item Z 1: Tätigkeiten gem. {\S}\,9
%     \item Z 2: Unterstützung des Arztes; Betreuung von Transporten
%     \item
%     Z 3: Verabreichung der \EE{Arzneimittelliste 1}:
%     Wird vom leitenden Arzt der Rettungsorganisation
%     festgelegt . Die Verabreichung kann
%     durchgeführt werden ohne dass ein Notarzt informiert werden muss oder es
%     ein solcher anordnet.
%     \item Z 4: eigenverantwortliche Betreuung von berufsspezifischen
%     Medizinprodukten und Arzneimitteln
%     \item 	Z 5: Mitarbeit in Forschung
%   \end{itemize}
%   \item[\T{NKA}] \marginlabel{{\S}\,11 Abs. 1 Z 1 SanG}
%   zusätzlich zu den oben genannten Kompetenzen darf der NKA jene
%   Medikamente der Arzneimittelliste 2 verabreichen, die ohne venösen
%   Zugang verabreicht werden können.
%   
%   Die \EE{Arzneimittelliste 2} wird vom leitenden
%   Arzt der Rettungsorganisation festgelegt und darf erst
%   verabreicht werden:
%   \begin{itemize}
%     \item auf Anweisung durch anwesenden Arzt (jeder mit jus
%     practicandi -- nicht unbedingt Notarzt) oder
%     \item nach Verständigung eines Notarztes.
%   \end{itemize}
%   \item[\T{NKV}] \marginlabel{{\S}\,11 Abs. 1 Z 2 SanG} zusätzlich zu den oben genannten
%   Kompetenzen darf der NKV jene Medikamenten der Arzneimittelliste 2 verabreichen, die über venösen Zugang verabreicht werden können.
%   \item[\T{NKI}] \marginlabel{{\S}\,12 SanG}
%      Durchführung einer endotrachealen Intubation samt Beatmung, aber ohne
%      Prämedikation
%      \begin{itemize}
%        \item
%        NKI-Kompetenz muss alle 2 Jahre rezertifiziert werden
%        \item
%        Durchführung der Intubation erst:
%        \begin{itemize}
%          \item
%          auf Anweisung durch anwesenden Arzt (jeder mit jus practicandi -- nicht unbedingt Notarzt) oder
%          \item
%          nach Verständigung eines Notarztes
%        \end{itemize}
%      \end{itemize}
% \end{labeling}

\subparagraph{Überschreiten der Kompetenzstufen}
Die Durchführung von Maßnahmen außerhalb des eigenen
Kompetenzstandes, führt zu einer Verwaltungsstrafe nach SanG bzw.
einer gerichtlichen Strafe wegen Körperverletzung und
zivilrechtlichen Schadenersatzanspruch. Selbst bei Wissen über Gefahren, Risiken und
Durchführungsart einer kompetenzmäßig höheren Maßnahme,
gibt es keine Rechtfertigung durch {\S}\,95 StGB (Unterlassene
Hilfeleistung).

Dies bedeutet weiters, dass einfache Maßnahmen, die zwar das
Berufsgesetz nicht aufzählt, aber von jedem Laien (mitunter nach
kurzer Einweisung) durchgeführt werden dürfen, auch von
Sanitätern gesetzt werden dürfen bzw. iSd. {\S}\,95 StGB gesetzt
werden müssen.
}{}{}{}{}

\Layout{2}{Rechtsfolge}{}{}{}{}{}

\Fazit{
Sepp darf keinen venösen Zugang legen oder sogar Medikamente darüber
verabreichen, selbst  wenn Gefahr in Verzug vorliegt. }


\subsubsection{Fall 8. Rechtswidrige Anordnung durch Notarzt}
% TODO : Subsubsection: Layout-Anpassung
\TODO{GABS}{yyyy-mm-dd}{Subsubsection: Layout-Anpassung notwendig.}{Dieser Unterunterabschnitt entspricht nicht dem Standard-Layout!}

\Layout{2}{Sachverhalt}{}{}{}{}{}

\begin{wrapfigure}{r}{5cm}
\includegraphics[width=\linewidth]{./images/teaser/emergency_physician_vehicle_00800.jpg}
\end{wrapfigure}Sepp konnte sich nun dazu überreden lassen keinen Venflon zu legen und
möchte wenigstens Nitro verabreichen. Gerade noch rechtzeitig bemerkt
er eine offene Packung von Viagra am Tisch liegen. Hr. Schmerz meint,
als er Sepp die Packung betrachten sieht, dass das
"`Zeug"' nichts nutzt, da er vor zwei
Stunden eine Tablette eingenommen hätte und noch immer
"`nichts steht"'. Kurze Zeit später
trifft der NEF mit der Notärztin Natascha ein, die den
Patienten flapsig untersucht und eine sofortige
Nitro-Verabreichung zur Bekämpfung des hohen Blutdrucks
(\RRmmHg{200}{110}) anordnet, obwohl sie von Sepp über die
gefundene Viagra-Packung informiert wurde.

\Layout{2}{Rechtsfrage}{Ist der Sanitäter verpflichtet eine rechtswidrige bzw. gefährliche
Anordnung des Notarztes (oder eines anderen Vorgesetzten) zu folgen?}{}{}{}{}



\Layout{2}{Rechtslage}{}{}{}{}{}

\EE{Ablehnung von rechtswidrigen Anordnungen}:

\begin{itemize}
    \item 
    Zivildiener: {\S}\,22 Abs. 2 Zivildienstgesetz:
    
    
    \begin{itemize}
        \item
        "`Er darf die Befolgung einer Weisung nur dann ablehnen,
        wenn [{\dots}] die Befolgung gegen strafgesetzliche Vorschriften
        verstoßen würde."'
        
    \end{itemize}
    
    \item 
    Haupt- und Ehrenamtliche: 
    
    \begin{itemize}
        \item 
        {\S}\,4 Abs. 1 S. 2 SanG
        
        \item 
        {\S}\,7 i.\,V.\,m. {\S}\,879 Abs. 1 ABGB
        
        \item 
        {\S}\,22 Abs. 2 ZDG analog:
        
    \end{itemize}
    
\end{itemize}

Aufgrund der im SanG fehlenden konkreten Bestimmung eines
Ablehnungsrechtes von Sanitätern, kann ein solches nur kompliziert
gewonnen werden aus:

\begin{itemize}
    \item allgemeinen Verpflichtung zur Wahrung des Patientenwohls nach {\S}\,4
    SanG    
    \item natürlichen Rechtsgrundsätzen iSd. {\S}\,7 ABGB und dem
    Nichtigkeitsgebot bei Sittenwidrigkeit des {\S}\,879 ABGB    
    \item analoger Anwendung des Zivildienstgesetz (ZDG) 
    \item Schutz der Grundrechte gem. Art.~2 EMRK bei
    Lebensgefahr und Art.~3 EGC bei sonstigen Verletzungen
    der Unversehrtheit.    
\end{itemize}

\Layout{2}{Rechtsfolge}{}{}{}{}{}

\Fazit{
Die RTW-Mannschaft darf daher die für den Patienten	schädigende Anordnung ablehnen. Natascha hingegen macht
sich strafbar wegen Körperverletzung und riskiert uU.
zivilrechtliche Schadenersatzansprüche
(${\rightarrow}$ Nitro-Gabe nicht lege artis).}


\subsubsection{Fall 9. Dokumentationspflicht}
% TODO : Subsubsection: Layout-Anpassung
\TODO{GABS}{yyyy-mm-dd}{Subsubsection: Layout-Anpassung notwendig.}{Dieser Unterunterabschnitt entspricht nicht dem Standard-Layout!}
\label{topic:dokumentation-rechtlich}

\Layout{2}{Sachverhalt}{Genervt von der Auseinandersetzung der Notärztin mit den Sanitätern,
beginnt Hr. Schmerz wild vor Schmerzen die Beteiligten anzubrüllen,
dass sie doch endlich ihre "`gottverdammte
Arbeit"' tun sollen. Der ebenso gereizte Bernd belehrt
daraufhin Hr. Schmerz, dass er sich doch gefälligst wie ein
"`richtiger Patient"' verhalten soll und
sich erst rühren darf, wenn es ihm gesagt wird. Diese Aussage wird
Bernd aber zum Verhängnis, da sich Hr. Schmerz Wochen später beim
Arbeitgeber von Alex, Sepp und Bernd über das
"`beleidigende Verhalten"' des Bernd
beschwert. Sepp seinerseits ist überzeugt, dass er sich korrekt
verhalten hat, doch kann er sich so recht an den alten Fall nicht mehr
erinnern.}{}{}{}{}



\Layout{2}{Rechtsfrage}{}{}{}{}{}

Besteht eine Dokumentationspflicht und wenn ja, in welchem Umfang?

\Layout{2}{Rechtslage}{\Z{{\S}\,5} \EE{SanG}:
"`\C{Sanitäter haben [{\dots}] die von
ihnen gesetzten Maßnahmen zu dokumentieren.}"'}{}{}{}{}~



\EE{Vertragsrechtliche Relevanz}:

Im österreichischen Recht hat grundsätzlich der Kläger den
Schaden, die Rechtswidrigkeit und das Verschulden des Beklagten zu
beweisen. Besteht aber zwischen Kläger und Beklagten ein Vertrag, so
wird gem. \Z{{\S}\,1298} ABGB das Verschulden des Beklagten
angenommen bzw. der Beklagte muss beweisen, dass ihn kein Verschulden trifft (=
\index{Beweislastumkehr}\T{Beweislastumkehr}).

Im Rettungswesen besteht zu jedem Patient ein \index{Behandlungsvertrag}\T{Behandlungsvertrag}, der eine
solche Beweislastumkehr bewirkt. Um nun für den Beklagten die
Möglichkeit zu schaffen, sein Unverschulden zu beweisen, wurde die
Dokumentationspflicht eingeführt, mit der bei ausreichender
Dokumentation wieder der Kläger das Verschulden des Beklagten zu
beweisen hat.


\EE{Inhalt der Dokumentation}:

\begin{itemize}
    \item gesetzte med. Maßnahmen
    (\prettyref{topic:dokumentation-medizinisch})
    \item situationsrelevante Umstände (beim Patient; am BO; während des
    Transportes; {\dots})
    \item Verweigerungen oder besondere Anmerkungen des Patienten	(kein Gurt; Marcoumarpflichtig; {\dots})
       \item uU. relevante Inhalte der Aufklärung
       \item sonstige für den Sanitäter entscheidend wirkende Umstände
\end{itemize}

% \Fazit{Grundsatz: Alles darf dokumentiert werden, es darf aber nicht Nichts
% dokumentiert werden.}

\Fazit{Grundsatz: Es muss gründlich und ausführlich
dokumentiert werden.}

\subparagraph{Durchführung der Dokumentation:}
Die Dokumentation soll vorrangig am Transportschein selbst oder wenn
nicht mehr ausreichend auf einem beigelegten Blatt Papier erfolgen.
Verweigert der Patient medizinische oder sonst notwendige
Maßnahmen, ist von ihm ein Revers zu unterschreiben, der dem
Transportschein beizulegen ist. Dabei hat der Sanitäter
immer (!) durch ein Gespräch auf mögliche Risiken und Gefahren der
Verweigerung hinzuweisen.

% Des Weiteren ist eine zusätzliche eigene Dokumentation für jeden
% Sanitäter empfehlenswert.

\Layout{2}{Rechtsfolge}{}{}{}{}{}

\Fazit{
Bernd hat keine ausreichende Dokumentation durchgeführt und kann sich
somit nicht gegen die Beschuldigungen zur Wehr setzen.}


\subsubsection{Fall 10. Verschwiegenheits- und Auskunftspflicht}
% TODO : Subsubsection: Layout-Anpassung
\TODO{GABS}{yyyy-mm-dd}{Subsubsection: Layout-Anpassung notwendig.}{Dieser Unterunterabschnitt entspricht nicht dem Standard-Layout!}

\Layout{2}{Sachverhalt}{Durch die Aufregung verschlechtert sich allmählich der
Zustand des Hr. Schmerz und das RTW-Team samt NEF-Mannschaft sind sich einig, dass ein
schleunigster Abtransport durchzuführen ist. Noch bevor sie aber den
Patienten umlagern können, erscheint Fr. Aufgeregt, die
Angehörige des Hr. Schmerz, und verlangt über alles informiert zu werden.}{}{}{}{}

\Layout{2}{Rechtsfrage}{Was umfasst die Verschwiegenheitspflicht der Sanitäter und Ärzte?
Besteht auch eine Auskunftspflicht derselben?}{}{}{}{}



\Layout{2}{Rechtslage}{}{}{}{}{}

Verschwiegenheitspflicht:

\begin{itemize}
    \item 
    \EE{strafrechtliche} Norm: \Z{{\S}~121} StGB
    \item 
    \EE{verwaltungsrechtliche} Norm: \Z{{\S}~6}~SanG
    \item 
    \E{beiden gleich}: Geheimnisse, die man aufgrund seiner
    Tätigkeit als Arzt/Sanitäter anvertraut bekommen hat, dürfen
    nicht an unbefugte Dritte weitergegeben werden
    \begin{itemize}
        \item \E{Geheimnisse}: alle Informationen, die
        einen Bezug zum Patienten haben oder liefern
        \item 
        daher darf man mit Kollegen oder sonstigen Personen über die erlebten
        Geschehnisse reden, solange aus dem Gesagten nicht erkennbar ist, um
        welche Person es sich bei dem Patienten handelt 
        \item 
        Bsp: Name des Patienten, Versicherungsnummer, Wohnadresse,
        {\dots}
    \end{itemize}
\end{itemize}

\EE{Wann \E{darf} man Geheimnisse weitergeben?}

\begin{itemize}
    \item Anonymität des Patienten bleibt gewahrt (=~kein
    Geheimnis)
    \item bei Entbindung von der Schweigepflicht durch den Patienten
\end{itemize}

\EE{Wann \E{muss} man Geheimnisse weitergeben?} = \T{Auskunftspflicht} \index{Auskunftspflicht} --
\Z{{\S}\,7} SanG

\begin{itemize}
    \item 
    wenn es zur Behandlung des Patienten dient
    \item 
    wenn es gesetzlichen vorgeschrieben ist (z.\,B. im
    Gerichtsverfahren als Zeuge)
    \item zum Schutz von höherwertigen Interessen der öffentlichen
    Gesundheitspflege oder Rechtspflege (${\rightarrow}$ zur Meldepflicht
    siehe unten)
\end{itemize}

\EE{Wem gegenüber muss man Auskunft erteilen?}

\begin{itemize}
    \item 
    Patienten
    \item 
    gesetzlichen Vertreter des Patienten
    \item 
    anderen Personen der Gesundheitsberufe, die den
    Patienten	betreuen, behandeln oder pflegen bzw. an denen der Patient im Zuge des Transportes übergeben wird
       \item Sozialversicherung und Krankenanstalt
       \item vom Patient als auskunftsberechtigt benannte Personen
\end{itemize}

\Layout{2}{Rechtsfolge}{}{}{}{}{}

\Fazit{
Fr. Aufgeregt zählt als bloße Angehörige nicht zum Kreis der
Auskunftsberechtigten, solange der Patient die Weitergabe von
Informationen nicht gestattet. }


\subsubsection{Fall 11. Fortbildungs- und Rezertifizierungspflicht}
\label{topic:fortbildungspflicht}
% TODO : Subsubsection: Layout-Anpassung
\TODO{GABS}{yyyy-mm-dd}{Subsubsection: Layout-Anpassung notwendig.}{Dieser Unterunterabschnitt entspricht nicht dem Standard-Layout!}
\label{topic:jus-fall-fortbildung}


\Layout{2}{Sachverhalt}{Nachdem nun auch Fr. Aufgeregt beruhigt wurde, wird Hr. Schmerz auf den
Sessel umgelagert, um über die enge Treppe den ersten Stock zu
verlassen. Doch handelt es sich beim Tragsessel um einen neuen
Stryker{\texttrademark}-Sessel, der dem Bernd noch vollkommen neu ist, da es den
Sessel zum Zeitpunkt seiner RS-Ausbildung noch gar nicht gegeben hat. Aber
genau wegen dieser Neuheit hatte sein Arbeitgeber eine interne
Fortbildung über die "`Neuen
Medizinprodukte"' abgehalten, zu der Bernd aber nicht
gekommen ist, da er meinte sowieso schon alles zu wissen und zu keinen
Fortbildungen verpflichtet ist. 

Sepp hat hingegen als vorbildlicher Praxisanleiter eine
Einschulung nach MPG erhalten.}{}{}{}{}



\Layout{2}{Rechtsfrage}{Besteht eine Fortbildungs- oder sogar Re\-zertifizierungs\-pflicht?}{}{}{}{}



\Layout{2}{Rechtslage}{}{}{}{}{}

\subparagraph{\EE{Fortbildungspflicht nach} \Z{{\S}~50} \EE{SanG}:}

Jeder Sanitäter ist dazu verpflichtet sich innerhalb von 2 Jahren ab
seinem Stichtag im Umfang von 16 Stunden fortzubilden. Als Stichtag
wird der Monatserste nach bestandener Rettungssanitäterausbildung
gewertet, wobei sich der Stichtag auch bei Erreichen von höheren
Kompetenzniveaus nicht ändert. Zu welchem Zeitpunkt die 16
Fortbildungsstunden innerhalb des Rahmens von 2 Jahren besucht werden,
ist dem Sanitäter überlassen.

Man kann also entweder regelmäßig Fortbildungen besuchen oder auf
einmal die Stundenzahl ansammeln. Weiters muss ein Sanitäter keine
bestimmten Themen als Fortbildungen besuchen, sondern kann nach seinem
Interesse vorgehen.

\subparagraph{\EE{Fortbildungspflicht nach} \Z{{\S}\,83}
\EE{MPG}:}

Die im Rettungswesen verwendeten Materialien gelten, abgesehen von
Arzneimitteln, als Medizinprodukte und dürfen daher iSd. MPG nur von
eingeschulten Personal angewendet werden. Diese Einschulungen müssen
nicht nur alle Medizinprodukte der Rettungsorganisation beinhalten,
sondern auch von allen Sanitätern besucht werden. Daher kann ein
Sanitäter in Bezug auf Stundenanzahl oder Fortbildungsthemen nicht
frei wählen. Die Fortbildungsstunden für das MPG können aber als
Fortbildungsstunden iSd. SanG gewertet werden. Zusätzlich ist zu
sagen, dass eine Einschulung der meisten Medizinprodukte schon durch
die Standardausbildung zum Rettungssanitäter abgedeckt ist. Für
die anderen, insbesondere höhere technische Produkte, muss jede
Organisation separate Fortbildungen anbieten.

\subparagraph{\EE{Rezertifizierungspflicht} \EE{gem}.\Z{{\S}~51} \EE{SanG}:}

Innerhalb des 2 Jahre Rahmens der Fortbildungspflicht muss sich ein
Sanitäter, unabhängig von seinem Ausbildungsstand, prüfen lassen
in der Herz-Lungen-Wiederbelebung mit halbautomatischen Defibrillator.
Der Rezertifizierungszeitpunkt innerhalb des 2 Jahre Rahmens kann vom
Sanitäter frei gewählt werden.

\E{Sonderfall NKI}: Ein Notfallsanitäter mit
Kompetenzen der Intubation muss sich zusätzlich innerhalb seines 2
Jahre Rahmens einmal über die Intubation prüfen lassen.

\Layout{2}{Rechtsfolge}{}{}{}{}{}

\Fazit{
Die für den vorliegenden Fall entscheidende Fortbildung gem. MPG wurde
von Bernd nicht besucht, was ihm nicht nur eine Verwaltungsstrafe,
sondern auch eine gerichtliche Strafe und zivilrechtlichen
Schadensersatzanspruch einbringen kann. Denn sollte es im
Umgang mit dem Medizinprodukt zu Körperverletzungen kommen, wird aufgrund der mangelnden
Materialkenntnis die Behandlung als nicht-lege-artis gewertet. Um einer
solchen Strafe zu entgehen genügt es, dass Bernd noch vor Ort, aber
vor Anwendung, durch eine "`geeignete
Person"' eingeschult wird.

Als "`geeignete Person"' gelten vorrangig
Medizinproduktebeauftragter,  Lehrsanitäter oder Praxisanleiter sowie
jene, die von der Rettungsorganisation als geeignet genannt werden.

Da Sepp Praxisanleiter ist, darf er vorab Bernd einweisen.
}


\subsection{Schadenersatz}
\label{topic:schadenersatz}

\subsubsection{Fall 12. Verletzung des Patienten durch
Sanitäter:}
% TODO : Subsubsection: Layout-Anpassung
\TODO{GABS}{yyyy-mm-dd}{Subsubsection: Layout-Anpassung notwendig.}{Dieser Unterunterabschnitt entspricht nicht dem Standard-Layout!}

\Layout{2}{Sachverhalt}{Nachdem nun Sepp unseren Bernd kurzum erklärt hat, wie er mit den
Stryker{\texttrademark}-Sessel umzugehen hat unterlässt er trotzdem das
Anschnallen des Patienten, da er meint, dass "`eh kana däs
mocht"'. Leider ist Hr. Schmerz kein Fliegengewicht
und der Treppenaufgang ist derart verwinkelt, dass es alle Mühe und
Not bedarf den Tragsessel ruhig zu halten.

Fast am Ende des Abgangs angekommen, flitzt auf einmal die
aufgeschreckte Katze des Hr. Schmerz an den beiden
"`Trägern"' vorbei, wodurch der
vorangehende Bernd seine Konzentration verliert, ausrutscht und samt
dem Patient seitlich nach vorne auf einen Kasten und die
darunter sich versteckende Katze namens
\emph{"`Kitty"'} %\emph{"`Muschi"'} 
fällt. Bernd kann sich
beim Sturz wenigsten selber abstützen und bleibt (auch durch die abfedernde Wirkung der Katze) unverletzt. Doch Hr. Schmerz prallt mit Schädel und Brustkorb
ungeschützt gegen den Kasten und wird unter dem hinter ihm
nachkommenden Tragsessel begraben. Er zieht sich neben einer RQW auf
der Stirn auch noch mehrere gebrochene Rippen zu.}{}{}{}{}


\Layout{2}{Rechtsfrage - Schadenersatz wegen Körperverletzung}{
Wann wird man schadenersatzpflichtig? Was umfasst der Schadenersatz? Wer
wird zur Haftung gezogen?
}
{}{}{}{}



\Layout{2}{Rechtslage}{}{}{}{}{}

\EE{Rechtliche Grundlage eines
Schadenersatzanspruches}:

\begin{itemize}
    
    \item allgemein: \Z{{\S}{\S}\,1293 ff ABGB}
    \item konkret für Körperverletzungen: \Z{{\S}{\S}
    1325-1327 ABGB}
\end{itemize}

\EE{Voraussetzungen der Schadenersatzpflicht}:

\begin{itemize}
    
    \item[1.] Schaden = Nachteil am Körper, an der Gesundheit, im Vermögen
    \item[2.] Kausalität = bei Wegdenken der Schädigerhandlung, bleibt auch der
    Erfolg aus
    \item[3.] Rechtswidrigkeit = objektiver Verstoß gegen Gesetze oder
    Vertragspflichten
    \item[4.] Schuld = subjektiver Verstoß gegen jene Sorgfaltspflichten, die
    man einhalten hätte können
\end{itemize}

\EE{Umfang des Schadenersatzes bei
Körperverletzungen}:

\begin{itemize}
	\item Heilungskosten
	\item Verdienstentgang
	\item Schmerzengeld
\end{itemize}


\EE{Wer haftet wie}:

Schädigt der Sanitäter im Zuge seines Dienstes einen Patient oder
dessen Angehörigen, kann der Geschädigte nach
\Z{{\S}\,1313a} ABGB direkt gegen die
Rettungsorganisation seine Schadenersatzansprüche geltend machen: Ein Sanitäter
gilt in dieser Beziehung als Erfüllungsgehilfe der
Rettungsorganisation, wodurch er als verlängerter Arm der
Organisation deren vertragliche Pflichten (sorgfältige Behandlung)
erfüllt. Daher kann der Geschädigte direkt gegen die Organisation
alle Ansprüche wie gegen den Sanitäter selber durchsetzen.

\Layout{2}{Rechtsfolge}{}{}{}{}{}

\Fazit{
Im vorliegenden Fall ist der Schaden durch RQW und Frakturen neben der
Kausalität gegeben. Weiters verstößt es gegen die gesetzlichen
und vertraglichen Pflichten den Patient ohne anzuschnallen zu
transportieren, weshalb eine Rechtswidrigkeit bejaht werden kann.

Auch hätten sich unsere Sanitäter anders, nämlich den Patient
anzuschnallen, verhalten können, wodurch ebenso die Schuld gegeben
ist. Daher kann Hr. Schmerz Schadenersatzansprüche gegen die Sanitäter oder,
wirtschaftlich vernünftiger, gegen die Rettungsorganisation geltend
machen.
}


\subsubsection{Fall 13. Sachbeschädigung}
% TODO : Subsubsection: Layout-Anpassung
\TODO{GABS}{yyyy-mm-dd}{Subsubsection: Layout-Anpassung notwendig.}{Dieser Unterunterabschnitt entspricht nicht dem Standard-Layout!}

\Layout{2}{Sachverhalt}{Leider sind die Verletzungen am Patient nicht die einzigen Schäden die
entstanden sind, denn der umgeschmissene Kasten war eine
Biedermeier-Rarität mit einem Wert von stolzen 10\,000,-- Euro. Auch die
Katze verwirkte als Auffangkissen für Bernd ihr Leben und hatte, neben
dem Sachwert von 350,-- Euro, auch einen besonderen emotionalen Wert für
Hrn. Schmerz.}{}{}{}{}



\paragraph{Rechtsfrage: Schadenersatz wegen Sachbeschädigung}

Wann wird man schadenersatzpflichtig? Was umfasst der Schadenersatz? Wer
wird zur Haftung gezogen?

\Layout{0}{Rechtslage}{

\subparagraph{Rechtliche Grundlage eines
Schadenersatzanspruches} Regelung im ABGB: \Z{\begin{inparaitem}
    \item Allgemein: \Z{{\S}{\S}\,1293 ff ABGB}; 
    \item Konkret für Sachschäden: \Z{{\S}\,1331 ABGB}
\end{inparaitem}}
% \begin{inparaitem}
%     \item Allgemein: \Z{{\S}{\S}\,1293 ff ABGB}; 
%     \item Konkret für Sachschäden: \Z{{\S}\,1331 ABGB}
% \end{inparaitem}

\subparagraph{Voraussetzungen der Schadenersatzpflicht}
Es gelten die gleichen Voraussetzungen wie bei Körperverletzungen.

\subparagraph{Umfang des Schadenersatzes bei Sachbeschädigung}
Maßgebend ist der \E{reale Schaden} (der tatsächlich eingetretener Schaden), ein 
\E{entgangener Gewinn} (ein in Zukunft drohender Schaden), sowie ein \E{ideeller Schaden}.
}{
\begin{itemize}
  \item Rechtliche Grundlage: ABGB
  \item Voraussetzung Schadenersatzpflicht:
  \begin{itemize}
    \item Wie bei Körperverletzungen
  \end{itemize}
  \item Umfang Schadenersatz bei Sachbeschädigung:
\begin{itemize}
    \item Realer Schaden (= tatsächlich eingetretener Schaden)
    \item Entgangener Gewinn (= in Zukunft drohender Schaden)
    \item Ideeller Schaden
  \end{itemize}
\end{itemize}
}{}{}{}

\Layout{2}{Rechtsfolge}{}{}{}{}{}

\Fazit{Gleiche Lösung wie bei Fall 12.}


\subsection{Behandlungspflicht}
\label{topic:behandlungspflichtJUS}

\subsubsection{Fall 14. Behandlungspflicht von Krankenanstalten, Ärzten,
Sanitätern und Durchschnittsbürgern:}
% TODO : Subsubsection: Layout-Anpassung
\TODO{GABS}{yyyy-mm-dd}{Subsubsection: Layout-Anpassung notwendig.}{Dieser Unterunterabschnitt entspricht nicht dem Standard-Layout!}

\Layout{2}{Sachverhalt}{
Vor Schmerzen und Trauer, um seine geliebte Katze
\E{"`Kitty"'},
%\E{"`Muschi"'},
beschimpft und verflucht Hr.
Schmerz die Anwesenden mit wilden Gestikulierungen. Als er dann
wenigstens von Alex und Sepp vorbildlich versorgt wird, beruhigt sich
Hr. Schmerz und lässt sich sogar mühelos ins Auto transportieren.

Nach ausreichender Versorgung des Hr. Schmerz begibt sich der
RTW auf die Reise zur Krankenanstalt. Doch ist leider das
einzig abbuchbare Bett, welches neben einer Internen- auch
eine Unfall-Ambulanz aufzuweisen hat, im Wilhelminenspital,
welches durch den großen Morgenstau gute 15 Minuten entfernt
ist. Schon nach kurzer Fahrzeit klagt Hr. Schmerz über
stärker werdende Atemprobleme und ein rasselndes, brodelndes
Atemgeräusch tönt durch den gesamten Patientenraum. Es hat
sich ein kardiales Lungenödem gebildet und Natascha fühlt
sich, trotz Verabreichung von Diuretika, nicht mehr in der
Lage den Patient während der Fahrt ausreichend zu behandeln.
Alex soll so schnell wie möglich das nächste Spital
anfahren. Durch eine kurze Abweichung der Route, gelangt das
Team zum AKH und prescht mit dem Patient zur nicht
vorinformierten Aufnahme.

Der dortige Aufnahmearzt Dr. Wasbringst hat schon einen 22 Stunden
Dienst hinter sich und ist alles andere als erfreut über das
"`Mitbringsel"' des Teams. Wutentbrannt
brüllt er Natascha an, was sie sich eigentliche erlaube trotz
gesperrten Bettenkontingent seine Station anzufahren und verweigert ihr
die Aufnahme des Hr. Schmerz.}{}{}{}{}



\Layout{2}{Rechtsfrage}{
Besteht eine Behandlungspflicht für Krankenanstalten? Besteht eine
Behandlungspflicht für medizinische Berufe? Besteht eine
ersthelferische Versorgungspflicht für den Durchschnittsbürger?
}{}{}{}{}


\Layout{0}{Rechtslage}{
\subparagraph{Behandlungspflicht für Krankenanstalten}
Jede Krankenanstalt, ob öffentlich oder privat, ist verpflichtet einen
behandlungspflichtigen Patienten  aufzunehmen und
zumindest ärztliche Grundbehandlung zukommen zu lassen.\ZFootnote{{\S}{\S}\,22 und 40 KAKuG}

Das in Wien geltende \E{Bettenkontingentsystem}, mit seinen Zuweisungen zu
bestimmten Krankenanstalten, ist trotzdem zulässig, da
sich der Patient während des Transportes schon in einem präklinischen Umfeld (=
Einsatzwagen mit Sanitätern und Medizinprodukten) befindet. Sollte
sich aber der Zustand des Patienten während der Fahrt
verschlechtern, darf das nächstgelegene Spital angefahren werden.
Erst wenn der Patient wieder transportfähig ist, darf er mit geeigneten
Krankentransportmitteln in eine andere Krankenanstalt weitergeschickt werden.

\subparagraph{Ärztliche Behandlungspflicht}
Ärzte, unabhängig ihrer Fachausbildung, sind dazu verpflichtet,
hilfsbedürftigen Patienten ärztliche Grundversorgung
zukommen zu lassen und -- wenn nötig -- weitere Einsatzkräfte nachzufordern oder
den Patienten in weitere medizinische Behandlung zu schicken.\ZFootnote{{\S}\,48 ÄrzteG}

\subparagraph{Behandlungspflicht für Sanitäter}
Jeder tätigkeits- oder berufsberechtigter Sanitäter, unabhängig
von seinem Kompetenzstand, ist zur selbstständigen und
eigenverantwortlichen Versorgung einer hilfsbedürftigen Person
verpflichtet. Dabei muss er, auch außerhalb des Dienstes, zumindest
qualifizierte Erste Hilfe leisten und wenn nötig einen Arzt (jeder
mit jus practicandi) nachfordern.\ZFootnote{{\S}\,4 i.\,V.\,m. {\S}\,8 i.\,V.\,m. {\S}\,9 Abs. 1 Z 1
SanG}


\subparagraph{Ersthelferische Versorgungspflicht für
Durchschnittsbürger}
Jeder Bürger (und nicht mehr tätigkeitsberechtigter Sanitäter) ist
zur Setzung von Hilfsleistungen verpflichtet, die seinem Wissens- und
Kenntnisstand entsprechen. Daher ist jeder mit einem
Führerschein-Erste-Hilfe-Kurs zur Setzung von zumutbaren Erste-Hilfe
Maßnahmen verpflichtet.\ZFootnote{{\S}\,95 StGB}

\subparagraph{Unzumutbarkeit der Behandlung}
Für alle Berufe oder Durchschnittsbürger ist die Behandlung dann
nicht zumutbar, wenn man \ldots
\begin{itemize}
  \item das eigene Leben oder die Gesundheit gefährden würde, oder
  \item höherwertige Interessen verletzen müsste.\ZFootnote{{\S}\,95 Abs. 2 StGB}
\end{itemize}
}{
\begin{itemize}
  \item Behandlungspflicht für Krankenanstalten. 
  \item Ärztliche Behandlungspflicht. 
  \item Behandlungspflicht für Sanitäter. 
  \item Ersthelferische Versorgungspflicht für
  Durchschnittsbürger. 
  \item Unzumutbarkeit der Behandlung. 
\end{itemize}
}{}{}{}    


\Layout{2}{Rechtsfolge}{
\Fazit{
Auch wenn das Bettenkontingent der Krankenanstalt aufgebraucht ist, muss Dr.
Wasbringst den Patient aufnehmen und zumindest soweit stabilisieren, dass
er ohne Gefahren in das eigentliche Zielspital weitertransportiert
werden kann.}
}{}{}{}{}




\subsection{Meldepflichtenrecht}

\subsubsection{Fall 15. Ansteckende Krankheiten}
% TODO : Subsubsection: Layout-Anpassung
\TODO{GABS}{yyyy-mm-dd}{Subsubsection: Layout-Anpassung notwendig.}{Dieser Unterunterabschnitt entspricht nicht dem Standard-Layout!}

\Layout{60}{Sachverhalt}{%
Als nun Hr. Schmerz nach langem hin und her auf der
Abteilung für \E{Innere Medizin} aufgenommen wurde, erholen
sich Alex, Sepp und Bernd bei einem kurzen Plausch mit der NEF-Mannschaft. Die Ruhe hält aber nicht lang an, als auch schon
der nächste Einsatz durchgegeben wird, der RTW und NEF,
zum selben BO schickt.

Dort angekommen finden sie in einer Lacke von Blut die
25-jährige Frau Hilflos liegen, die sich selber (querverlaufend) die Pulsadern
aufgeschnitten hat. Als Sepp beginnt sie zu versorgen, erzählt Fr.
Hilflos, dass ihr Leben "`sowieso keinen Sinn mehr
hat"', da sie ihr Freund mit \index{HIV!Meldepflicht}\Abk{HIV} angesteckt hat.
}
{}
{./images/teaser/blut-tatort-jacke.jpg}
{}
{GABS.}





\Layout{2}{Rechtsfrage}{
Bei welchen Erkrankungen besteht eine Meldepflicht? In welchem
Ausmaß besteht eine solche Meldepflicht?
}{}{}{}{}



\Layout{2}{Rechtslage}{}{}{}{}{}


\EE{Meldepflichtige Krankheiten}:

\begin{itemize}
  \item
  grundsätzlich:
  \begin{itemize}
    \item
    leicht übertragbare Krankheiten
    \item
    mit Gefahr von schweren Gesundheits- oder Todesfolgen
    \item
    Meldepflicht nur wenn es Gesetz oder Verordnung anordnen
  \end{itemize}
  
  \item
  nach \E{Epidemiegesetz}:
  
  \begin{itemize}
    \item
    bei bloßen Verdacht oder sicherer Diagnose einer im Gesetz genannten
    Erkrankung.
  \end{itemize}
  \item
  nach \E{Tuberkulosegesetz}:
  \begin{itemize}
    \item
    bei bloßen Verdacht auf eine Erkrankung die mit Sicherheit oder
    Wahrscheinlichkeit durch Tuberkulosebakterien verursacht wurde
  \end{itemize}
  \item
  nach \E{Geschlechtskrankheitengesetz}:
  \begin{itemize}
    \item
    nur bei Befürchtung, dass sich die Krankheit	weiterverbreiten würde oder sich der Patient einer
       Heilbehandlung entziehen würde.
  \end{itemize}
  \item
  nach \E{AIDS-Gesetz}:
  \begin{itemize}
    \item
    \E{Anonymisierte} Meldepflicht erst ab: manifester
    Erkrankung mit AIDS (= bei sichtbaren Krankheitszeichen;
    nicht schon bei HIV-Infektion) oder Todesfall mit
    erkennbaren AIDS typischen Krankheiten\footnote{s.
    \prettyref{topic:aids}}
    \item
    es gibt keine Untersuchungs- oder Behandlungspflicht daher darf kein
    behördlicher Zwang ausgeübt werden
  \end{itemize}
  \item
  nach \E{Verordnungen} des Gesundheitsministerium:
  \begin{itemize}
    \item da sich ständig neue übertragbare
    Krankheiten\footnote{in letzter Zeit z.\,B.
    \E{SARS}, \E{Vogelgrippe}, \E{Schweinegrippe}
    (\prettyref{topic:grippe})} entwickeln, gibt das
    Gesundheitsministerium in unregelmäßigen  Abständen Verordnungen heraus, die eine Meldepflicht beinhalten.
  \end{itemize}
\end{itemize}

\EE{Wer ist zur Meldung verpflichtet?}

Nicht jeder Durchschnittsbürger ist zur Meldung verpflichtet, sondern
nur besondere Personengruppen worunter aber alle Gesundheitsberufe
fallen.

\EE{An wen ist zu melden?}

\begin{itemize}
    
    \item 
    Bei AIDS: Gesundheitsministerium
    
    \item 
    Ansonsten: Bezirksverwaltungsbehörde bzw. in Wien an MA 15
    
\end{itemize}

\Layout{2}{Rechtsfolge}{}{}{}{}{}

\Fazit{%
Nun ist zwar unser RTW-Team nicht zur Meldung von HIV an
eine Behörde verpflichtet, doch sehr wohl zur Meldung an die weiterbehandelnden
Gesundheitsberufe aufgrund der schon erwähnten Auskunftspflicht gem.
\Z{{\S}\,7} SanG.}


\subsubsection{Fall 16. Strafbare Handlungen}
% TODO : Subsubsection: Layout-Anpassung
\TODO{GABS}{yyyy-mm-dd}{Subsubsection: Layout-Anpassung notwendig.}{Dieser Unterunterabschnitt entspricht nicht dem Standard-Layout!}
\label{topic:strafbare-handlungen}


\Layout{60}{Sachverhalt}{%
Die Schnittwunden am Handgelenk sind nicht die einzigen Verletzungen der
Fr. Hilflos. 

Mehrere Hämatome\footnote{\emph{Hämatom:}
Bluterguss.} überziehen den Körper und ein großes
Veilchen ziert das rechte Auge der jungen Patienten. Als Natascha auf diese weiteren Verletzungen
eingeht, beginnt Fr. Hilflos weinend zu berichten, dass ihr Freund sie öfters geschlagen hat und er ihr mit
dem Tod gedroht hat, falls sie die Polizei verständigen würde.
}
{}
{./images/trauma/haematom-fehlpunktion-001.jpg}
{Hämatom.}
{GABS.}



\Layout{2}{Rechtsfrage}{
Besteht eine Pflicht zur Anzeige von strafbaren Handlungen? Wenn ja,
für wen besteht diese Anzeigepflicht?
}{}{}{}{}



\Layout{2}{Rechtslage}{}{}{}{}{}


\subparagraph{\T{Anzeigepflicht}}
\label{topic:anzeigepflicht}
\index{Anzeigepflicht}
Nur bestimmte Personengruppen sind dazu verpflichtet bei Kenntnis über
strafbare Handlungen dies der Polizei oder Staatsanwaltschaft per
Anzeige zu melden.

\begin{itemize}
  \item
  Ärztliche Anzeigepflicht\ZFootnote{{\S}\,54 Abs. 4 und 5
  ÄrzteG}:
  
  Nur bei Vorliegen folgender Delikte (Verdacht reicht aus) und
  unabhängig vom Willen des Betroffenen:
  \begin{itemize}
    \item
    Tötung bzw. schwere Körperverletzung
    \item
    Misshandlung, Quälen, Vernachlässigung oder sexueller Missbrauch
    von:
    \begin{itemize}
      \item Minderjährigen oder
      \item Volljährigen, die nicht alleine in der Lage sind ihre Interessen
      wahrzunehmen
    \end{itemize}
  \end{itemize}
  \item Polizeiliche Ermittlungspflicht\ZFootnote{{\S}\,99 StPO}
  für alle Delikte
\end{itemize}


\subparagraph{\T{Anzeigerecht}}
\label{topic:anzeigerecht}
\index{Anzeigerecht}

Jeder Durchschnittsbürger und auch sonstige Gesundheitsberufe sind
\E{nicht verpflichtet} zur Anzeige von Strafdelikten. Sie haben aber das
\E{Recht}, eine strafbare Handlung anzuzeigen.\ZFootnote{{\S}\,80 Abs. 1 StPO}

Trotz der fehlenden Verpflichtung für den Sanitäter sollte bei jedem
Einsatz, bei dem der Verdacht einer strafbaren Handlung vorliegt,
zumindest die Polizei nachgefordert werden. Die Nachalarmierung kann
man in weiter Analogie mit der aus \Z{{\S}\,4} SanG
entstammenden Pflicht für den Sanitäter "`das Wohl der
Patienten [{\dots}] zu wahren"' begründen.

\Layout{2}{Rechtsfolge}{
\Fazit{%
Natascha muss aufgrund ihrer gesetzlichen Anzeigepflicht die Polizei
nachfordern, die zur Einbringung einer Anzeige den Sachverhalt
aufnehmen wird. Dies muss sogar dann geschehen, wenn die
Patientin selber keine Anzeige einbringen möchte.}
}{}{}{}{}




\subsection{Straßenordnung}
\label{topic:strassenverkehrsordnung}

\subsubsection{Fall 17. Blaulichtverwendung}
% TODO : Subsubsection: Layout-Anpassung
\TODO{GABS}{yyyy-mm-dd}{Subsubsection: Layout-Anpassung notwendig.}{Dieser Unterunterabschnitt entspricht nicht dem Standard-Layout!}

\Layout{60}{Sachverhalt}{%
Nach Durchführung der Behandlung und Anzeige gegen den aggressiven
Freund, begibt sich das RTW-Team erschöpft zum Stützpunkt zurück.
Hungrig und müde quält sich die RTW-Mannschaft durch das
alltägliche Wiener Verkehrschaos, als Sepp auf einmal meint:
\Speech{"`Geh, schalt doch das Hörndl und die Blitzer ei,
damit dee uns endli Plotz machen und wir einrücken
kännan!"'}
}
{}
{./images/gabriel-sebastian-aass/IMG_1896_00800.jpg}
{}
{GABS.}



\Layout{2}{Rechtsfrage}{%
Wann darf das Blaulicht samt Folgetonhorn verwendet werden? Wie darf
bzw. muss man sich bei einer Blaulichtfahrt verhalten?
}{}{}{}{}



\Layout{2}{Rechtslage}{}{}{}{}{}

\EE{Rechtsquellen}: \Z{{\S}\,26}~StVO und \Z{{\S}\,107}~KFG

\subparagraph{Verwendung von Blaulicht und Folgetonhorn:}

\begin{itemize}
    \item Grundsätzlich bei \E{Gefahr in Verzug}; im Näheren bei:
    \begin{itemize}
        \item Fahrt zum Berufungsort
        \item Haltezeit am Berufungsort
        \item Fahrt zum Abgabeort
    \end{itemize}
\end{itemize}

\subparagraph{Verhalten bei Blaulichtfahrt:}

\begin{itemize}
    \item Verkehrsgebote und -beschränkungen gelten nicht
    \item Geschwindigkeitshöchstgrenzen gelten nicht
    \item Ampeln gelten als Stopptafeln
    \item Absoluter Vorrang vor Nicht-Einsatzfahrzeugen
    \item Es dürfen weder Personen noch Sachen gefährdet oder beschädigt
    werden
    \item Es darf kein Nicht-Einsatzfahrzeug dem Rettungswagen unmittelbar
    nachfahren
    \item Vertrauensgrundsatz ist ungültig (${\rightarrow}$ man darf bei
    Blaulichtfahrten nicht mehr darauf vertrauen, dass sich die anderen
    Verkehrsteilnehmer richtig verhalten; daher ist defensives
    und voraussehendes Fahren verpflichtend)
\end{itemize}

\subparagraph{Vorrang der Einsatzfahrzeuge untereinander:}

\begin{enumerate}
    \item Rettungsfahrzeuge
    \item Feuerwehrfahrzeuge
    \item Polizeifahrzeuge
    \item sonstige Einsatzfahrzeuge
\end{enumerate}

\subparagraph{Missbrauch von Blaulicht und Folgetonhorn:}


Wer Blaulicht oder Folgetonhorn in Situationen verwendet, bei denen
nicht Gefahr in Verzug vorliegt, macht sich gerichtlich strafbar nach
\Z{{\S}\,1} des \E{Bundesgesetzes gegen
den Missbrauch von Notzeichen}.

\subparagraph{Gebot bei Blaulichtfahrten}
\Fazit{%
Nur wer noch größere Sorgfalt und Vorsicht wahrt als bei normaler Fahrt, verhält sich richtig!

\EE{Das beinhaltet mitunter eine \underline{geringere} Fahrgeschwindigkeit als
die höchst zulässige!} }

\subparagraph{Exkurs: Anschnallpflicht}

Der Fahrer eines jeden Fahrzeuges ist dazu verpflichtet zu achten, dass
er und seine Passagiere angeschnallt sind.\ZFootnote{{\S}\,106~KFG} Es steht aber jedem
Passagier frei zu sagen, dass er unangeschnallt mitfahren möchte,
wodurch eine Schuld beim Fahrer entfällt. Für Kranken- und
Rettungstransport besteht zusätzlich gem. Kraftfahrzeuggesetz (KFG) die
Ausnahme, dass ein Passagier (Personal wie Patient) unangeschnallt mitfahren darf, wenn es ansonsten
"`mit dem Zweck der Fahrt unvereinbar
ist"'.\ZFootnote{{\S}\,106 Abs.\,3~Z\,3 KFG} Trotzdem sollte
der Patient aufgefordert werden sich anzuschnallen bzw. sollte vom Sanitäter angeschnallt werden.
Verweigert der Patient, muss diese Verweigerung dokumentiert werden (am
besten mittels Revers, den der Patient unterschreibt).

Wird man nun im Zuge eines (sogar unverschuldeten) Verkehrsunfalls
verletzt und war man dabei nicht angeschnallt, verliert man einen Teil
des Schandenersatzanspruchs, weil das Fahren ohne angelegten Gurt als
Mitverschulden gewertet wird.

\Layout{2}{Rechtsfolge}{}{}{}{}{}

\Fazit{%
Würde unser RTW-Team tatsächlich Blaulicht und Folgetonhorn
einschalten, nur um schneller einrücken zu können, würden sie sich
für die falsche Blaulichtfahrt gerichtlich strafbar machen. Bei einem
-- auch  unverschuldeten -- Unfall müssten unsere drei Sanitäter die
volle Haftung tragen.}


\vspace{1em}
Der weitere Tag verläuft ruhig mit den üblichen Patienten
bis zum \ldots

\vspace{1em}
{\centering
\EE{{\sffamily\Large Feierabend.}}
\par}


